\documentclass[12pt]{article}
\usepackage[a4paper,margin=1in]{geometry}
\usepackage{amsmath,amssymb}
\usepackage{enumitem}
\usepackage{mathtools}
\usepackage{hyperref}
\usepackage{titlesec}
\usepackage{amsthm}

\title{Group \& Ring Theory MATP33: Course Summary of Important Results}
\author{Simon Gustafsson}
\date{}

\newtheorem{theorem}{Theorem}
\newtheorem{lemma}{Lemma}
\newtheorem{corollary}{Corollary}
\newtheorem{definition}{Definition}
\newtheorem{remark}{Remark}


\begin{document}
\maketitle

\section*{Lecture 0 — Background Results}

\begin{theorem}[First Isomorphism Theorem]
Let $f : G \to H$ be a group homomorphism. Then $\ker f \trianglelefteq G$ and $\operatorname{im} f \cong G / \ker f$.
\end{theorem}

\begin{theorem}[Second Isomorphism Theorem]
Let $G$ be a group, $H \leq G$, and $N \trianglelefteq G$. Then $HN \leq G$, $H \cap N \trianglelefteq H$, and $HN/N \cong H/(H \cap N)$.
\end{theorem}

\begin{theorem}[Third Isomorphism Theorem]
If $K \trianglelefteq G$ and $K \leq L \trianglelefteq G$, then $(G/K)/(L/K) \cong G/L$.
\end{theorem}

\begin{theorem}[Correspondence Theorem]
Let $N \trianglelefteq G$. Every subgroup of $G/N$ is of the form $H/N$ where $N \leq H \leq G$, and conversely if $N \leq H \leq G$ then $H/N \leq G/N$.
\end{theorem}

\begin{theorem}[Lagrange]
If $H \leq G$ and $G$ is finite, then $|H|$ divides $|G|$.
\end{theorem}

\begin{theorem}[Fundamental Theorem of Finitely Generated Abelian Groups]
Let $A$ be a finitely generated abelian group. Then
\[
  A \cong \mathbb{Z}^s \oplus \mathbb{Z}/a_1\mathbb{Z} \oplus \cdots \oplus \mathbb{Z}/a_r\mathbb{Z},
\]
where $s \geq 0$ and $a_1 \mid a_2 \mid \cdots \mid a_r$ are non-zero non-units in $\mathbb{Z}$. This decomposition is unique.
\end{theorem}

\section*{Lecture 1}

\begin{theorem}[1.10]
Let $G$ be a group and $X$ a set. Then:
\begin{enumerate}[label=(\roman*)]
  \item If $X$ is a $G$-set, then the action of $G$ on $X$ induces a homomorphism $\varphi : G \to S_X$;
  \item Any homomorphism $\varphi : G \to S_X$ induces an action of $G$ on $X$.
\end{enumerate}
\end{theorem}

\begin{proof}
(i) Define $\varphi : G \to S_X$ by $(\varphi(a))(x) = a \cdot x$ for $a \in G$ and $x \in X$.
We must first check that $\varphi(a) \in S_X$ for each $a \in G$, i.e.\ that the function $x \mapsto a \cdot x$ is a bijection $X \to X$.
\begin{itemize}
  \item \textit{Injective:} Suppose $a \cdot x = a \cdot y$. Then $x = a^{-1} \cdot (a \cdot x) = a^{-1} \cdot (a \cdot y) = y$.
  \item \textit{Surjective:} For any $x \in X$, we have $a \cdot (a^{-1} \cdot x) = x$, so $x \in \operatorname{im}(\varphi(a))$.
\end{itemize}
Hence $\varphi(a) \in S_X$. To see $\varphi$ is a homomorphism, let $a, b \in G$. Then for all $x \in X$,
\[
  (\varphi(ab))(x) = (ab) \cdot x = a \cdot (b \cdot x) = (\varphi(a))((\varphi(b))(x)) = (\varphi(a)\varphi(b))(x).
\]

(ii) Define $a \cdot x := (\varphi(a))(x)$. Then $(ab) \cdot x = (\varphi(ab))(x) = (\varphi(a)\varphi(b))(x) = a \cdot (b \cdot x)$, and $e \cdot x = (\varphi(e))(x) = x$. Hence $X$ is a $G$-set.
\end{proof}


\begin{theorem}[1.11 — Cayley's Theorem]
Every group $G$ is isomorphic to a subgroup of the symmetric group $S_G$.
\end{theorem}

\begin{proof}
By Theorem 1.10(i), left translation $a \cdot x := ax$ gives a homomorphism $\varphi : G \to S_G$.
If $a \in \ker\varphi$, then $\varphi(a) = \mathrm{id}_G$, so $ax = x$ for all $x \in G$, giving $a = e$.
Hence $\ker\varphi = \{e\}$, and by the First Isomorphism Theorem $G \cong G/\ker\varphi \cong \operatorname{im}\varphi \leq S_G$.
\end{proof}


\begin{definition}[Normal core]
Let $G$ be a group and $H \leq G$. The \emph{normal core} of $H$ in $G$ is
\[
  \operatorname{core}(H) := \bigcap_{x \in G} xHx^{-1}.
\]
It is the largest normal subgroup of $G$ contained in $H$.
\end{definition}

\begin{theorem}[1.12]
Let $G$ be a group and $H \leq G$ a subgroup of index $n$. Then there is a homomorphism $\varphi : G \to S_n$ such that $\ker\varphi = \operatorname{core}(H)$.
\end{theorem}

\begin{proof}
Since $|G/H| = n$, we have $S_{G/H} \cong S_n$. By Theorem 1.10(i), the action of $G$ on the coset space $G/H$ by left multiplication gives a homomorphism $\varphi : G \to S_n$ with $\varphi(a)(xH) = axH$. Then
\begin{align*}
  \ker\varphi &= \{a \in G \mid axH = xH \text{ for all } x \in G\} \\
              &= \{a \in G \mid x^{-1}ax \in H \text{ for all } x \in G\} \\
              &= \{a \in G \mid a \in xHx^{-1} \text{ for all } x \in G\} \\
              &= \bigcap_{x \in G} xHx^{-1} = \operatorname{core}(H). \qedhere
\end{align*}
\end{proof}

\begin{corollary}[1.13]
Let $G$ be a group with a normal subgroup $H$ of index $n$. Then $G/H$ is isomorphic to a subgroup of $S_n$.
\end{corollary}

\begin{proof}
Since $H \trianglelefteq G$, conjugation fixes $H$, so $\operatorname{core}(H) = \bigcap_{x \in G} xHx^{-1} = H$. By Theorem 1.12, $\ker\varphi = H$, and by the First Isomorphism Theorem $G/H \cong \operatorname{im}\varphi \leq S_n$.
\end{proof}


\begin{theorem}[1.24 — Orbit-Stabiliser Theorem]
Let $G$ be a group acting on a set $X$. Then the set of all orbits in $X$ under $G$ is a partition of $X$. Furthermore, for each $x \in X$ there is a bijection $Gx \to G/G_x$, and hence
\[
  |Gx| = |G : G_x|.
\]
Therefore, if $X$ is finite,
\[
  |X| = \sum_{x \in C} |G : G_x|,
\]
where $C \subseteq X$ contains exactly one element from each orbit.
\end{theorem}

\begin{proof}
\textit{(Partition.)} Define $x \sim y$ if $x = ay$ for some $a \in G$. This is an equivalence relation: reflexivity follows from $e \cdot x = x$, symmetry from $x = ay \Rightarrow y = a^{-1}x$, and transitivity from $x = ay,\, y = bz \Rightarrow x = (ab)z$. The equivalence class of $x$ is exactly the orbit $Gx$, so the orbits partition $X$.

\textit{(Bijection $Gx \to G/G_x$.)} Define
\[
  \varphi : Gx \to G/G_x, \qquad ax \mapsto aG_x.
\]
The key observation is the following chain of equivalences, for any $a, b \in G$:
\[
  ax = bx \iff a^{-1}bx = x \iff a^{-1}b \in G_x \iff aG_x = bG_x.
\]
\begin{itemize}
  \item \textit{Well-defined:} If $ax = bx$ (the same orbit element written using two different group elements), then $aG_x = bG_x$, so the coset $\varphi$ assigns does not depend on the choice of representative.
  \item \textit{Injective:} If $\varphi(ax) = \varphi(bx)$, i.e.\ $aG_x = bG_x$, then $ax = bx$ by the chain above.
  \item \textit{Surjective:} Take any coset $aG_x \in G/G_x$. Since $a \in G$, we have $ax \in Gx$, and $\varphi(ax) = aG_x$.
\end{itemize}
Hence $\varphi$ is a bijection, giving $|Gx| = |G/G_x| = |G : G_x|$. Summing over one representative per orbit gives $|X| = \sum_{x \in C} |G : G_x|$.
\end{proof}


\begin{theorem}[1.25 — Class Equation]
Let $G$ be a finite group. Then the conjugacy classes partition $G$, and
\[
  |G| = \sum_{a \in C} |G : C_G(a)|,
\]
where $C \subseteq G$ contains exactly one element from each conjugacy class.
\end{theorem}

\begin{proof}
Let $G$ act on itself by conjugation, i.e.\ $a \cdot x = axa^{-1}$. The orbit of $x$ under this action is its conjugacy class $\operatorname{ccl}(x) = \{axa^{-1} \mid a \in G\}$, and the stabiliser of $x$ is $G_x = \{a \in G \mid axa^{-1} = x\} = C_G(x)$. Applying Theorem 1.24 directly gives the partition into orbits and $|G| = \sum_{a \in C} |G : C_G(a)|$.
\end{proof}


\section*{Lecture 2}

\begin{theorem}[2.1]
Let $G$ be a finite group of order $p^n$, where $p$ is prime and $n \in \mathbb{N}$. Then:
\begin{enumerate}[label=(\roman*)]
  \item $G$ has a non-trivial centre $Z(G)$;
  \item $Z(G) \cap K$ is non-trivial for any non-trivial normal subgroup $K$ of $G$;
  \item if $H$ is a proper subgroup of $G$, then $H \subsetneq N_G(H)$. In particular, if $|H| = p^{n-1}$ then $H \trianglelefteq G$.
\end{enumerate}
\end{theorem}

\begin{proof}
(i) Since $x \in Z(G)$ iff $\operatorname{ccl}(x) = \{x\}$, the class equation refines to
\[
  p^n = |Z(G)| + \sum_{x \in C} |G : C_G(x)|,
\]
where $C$ runs over conjugacy classes with more than one element. For each $x \in C$, $|C_G(x)|$ properly divides $p^n$, so $p \mid |G : C_G(x)|$. Hence $p$ divides the sum, and therefore $p \mid |Z(G)|$. Since $e \in Z(G)$, we have $|Z(G)| \geq 1$, so in fact $|Z(G)| \geq p > 1$.

(ii) Intersecting the class decomposition of $G$ with $K$ gives $|K| = |Z(G) \cap K| + \sum_{x \in C} |\operatorname{ccl}(x) \cap K|$. Since $K \trianglelefteq G$, each conjugacy class is either disjoint from $K$ or contained in $K$, so $p \mid \sum |\operatorname{ccl}(x) \cap K|$. As $p \mid |K|$, we get $p \mid |Z(G) \cap K|$, and the same argument as (i) gives $Z(G) \cap K$ non-trivial.

(iii) Let $M$ be a maximal normal subgroup of $G$ contained in $H$. Then $G/M$ is a p-group, so by (i) it has a non-trivial centre $L/M \trianglelefteq G/M$. By the Correspondence Theorem, $L \trianglelefteq G$. Since $M \subsetneq L$ strictly, maximality of $M$ forces $L \not\subseteq H$.

It remains to show $L \subseteq N_G(H)$. To show $\ell \in N_G(H)$ we need $\ell^{-1}h\ell \in H$ for all $h \in H$. Take $h \in H$ and $\ell \in L$. Since $\ell M \in L/M = Z(G/M)$, it commutes with every coset, so $(\ell M)(hM) = (hM)(\ell M)$, i.e.\ $\ell h M = h\ell M$. This gives $\ell^{-1}h\ell \in hM$. Since $h \in H$ and $M \subseteq H$, we have $hM \subseteq H$, so $\ell^{-1}h\ell \in H$. Hence $\ell \in N_G(H)$, so $L \subseteq N_G(H)$.

Since $L \subseteq N_G(H)$ but $L \not\subseteq H$, we have $H \subsetneq N_G(H)$. Finally, if $|H| = p^{n-1}$ then $N_G(H)$ has order strictly greater than $p^{n-1}$, and the only possibility is $p^n = |G|$, so $N_G(H) = G$ and $H \trianglelefteq G$.
\end{proof}


\begin{theorem}[2.2 — Burnside's Lemma]
Let $G$ be a finite group acting on a finite set $X$, and for $g \in G$ let $X_g := \{x \in X \mid gx = x\}$. Then the number $k$ of orbits in $X$ under $G$ is
\[
  k = \frac{1}{|G|} \sum_{g \in G} |X_g|.
\]
\end{theorem}

\begin{proof}
Let $S := \{(g, x) \in G \times X \mid gx = x\}$. Fixing $g$ gives $|S| = \sum_{g \in G}|X_g|$; fixing $x$ gives $|S| = \sum_{x \in X}|G_x|$. Let $C \subseteq X$ contain one element from each orbit. By the Orbit-Stabiliser Theorem, $|G_x| = |G|/|Gx| = |G|/|Ga|$ for all $x \in Ga$. Hence
\[
  \frac{1}{|G|}\sum_{g \in G}|X_g|
  = \frac{|S|}{|G|}
  = \frac{1}{|G|}\sum_{x \in X}|G_x|
  = \frac{1}{|G|}\sum_{a \in C}\sum_{x \in Ga}\frac{|G|}{|Ga|}
  = \frac{1}{|G|}\sum_{a \in C}|G|
  = k. \qedhere
\]
\end{proof}


\begin{remark}[Pólya counting]
Burnside's Lemma counts distinct objects under symmetry. As an illustration: how many necklaces with $p$ beads ($p$ prime) can be made from $n$ colours, where two necklaces are the same if one is a rotation of the other?

Let $X$ be the set of all colourings, so $|X| = n^p$. The cyclic group $G = \langle a \rangle$ of order $p$ acts on $X$ by rotation. For the identity, $X_e = X$ so $|X_e| = n^p$. For any $g \neq e$: since $|G|$ is prime, $g$ generates $G$, so a necklace fixed by $g$ must be fixed by all rotations, meaning all beads have the same colour. Hence $|X_g| = n$. Burnside then gives
\[
  k = \frac{1}{p}\left(n^p + (p-1)n\right).
\]
\end{remark}


\section*{Lecture 3}

\begin{definition}[p-group, Sylow p-subgroup]
Let $p$ be a prime. A group $G$ is a \emph{$p$-group} if every element of $G$ has order a power of $p$. If $G$ is finite and $p^m \mid |G|$ but $p^{m+1} \nmid |G|$, then any subgroup of $G$ of order $p^m$ is called a \emph{Sylow $p$-subgroup} of $G$.
\end{definition}

\begin{theorem}[3.5 — First Sylow Theorem]
Let $G$ be a finite group and $p$ a prime. If $p^m \mid |G|$, then $G$ has a subgroup of order $p^m$.
\end{theorem}

\begin{proof}
If $m = 0$ the result is trivial. We proceed by induction on $n = |G|$, the case $n = 1$ being trivial. Assume the result holds for all groups of order less than $n$.

\textit{Case 1: $p \mid |Z(G)|$.} By the Fundamental Theorem of Finitely Generated Abelian Groups, $Z(G)$ contains an element $a$ of order $p$. Let $C = \langle a \rangle$. Since $a \in Z(G)$, we have $C \trianglelefteq G$, and $|G/C| = n/p$, which is divisible by $p^{m-1}$. By the induction hypothesis, $G/C$ has a subgroup of order $p^{m-1}$, which by the Correspondence Theorem is of the form $H/C$ for some $H \leq G$. Then $|H| = |H/C| \cdot |C| = p^{m-1} \cdot p = p^m$.

\textit{Case 2: $p \nmid |Z(G)|$.} From the class equation $n = |Z(G)| + \sum_a |G : C_G(a)|$, since $p \mid n$ but $p \nmid |Z(G)|$, there must exist some $a \notin Z(G)$ with $p \nmid |G : C_G(a)|$. Since $|G| = |G : C_G(a)| \cdot |C_G(a)|$ and $p^m \mid |G|$ but $p \nmid |G : C_G(a)|$, it follows that $p^m \mid |C_G(a)|$. Since $a \notin Z(G)$, we have $C_G(a) \neq G$, so $|C_G(a)| < n$. By the induction hypothesis $C_G(a)$ contains a subgroup of order $p^m$, which is also a subgroup of $G$.
\end{proof}

\begin{corollary}[3.6 — Cauchy's Theorem]
If $p$ is a prime dividing $|G|$, then $G$ has an element of order $p$.
\end{corollary}

\begin{proof}
By the First Sylow Theorem, $G$ has a subgroup $H$ of order $p$. Every non-identity element of $H$ has order $p$.
\end{proof}

\begin{corollary}[3.7]
A finite group $G$ is a $p$-group if and only if $|G|$ is a power of $p$.
\end{corollary}

\begin{proof}
If $|G|$ is a power of $p$, then by Lagrange every element order divides $|G|$, hence is a power of $p$. Conversely, if $|G|$ is not a power of $p$, then some prime $q \neq p$ divides $|G|$, and by Cauchy's Theorem $G$ has an element of order $q$, contradicting that $G$ is a $p$-group.
\end{proof}



\begin{theorem}[3.8 — Second Sylow Theorem]
Let $G$ be a finite group and $p$ a prime. All Sylow $p$-subgroups of $G$ are conjugate.
\end{theorem}

\begin{proof}
Let $S, H \in \mathrm{Syl}_p(G)$. Let $S$ act on the coset space $G/H$ by left translation: $s \cdot gH = sgH$.

\textit{Orbits have $p$-power size.} By the Orbit-Stabiliser Theorem, each orbit has size $|S : S_{gH}|$, which divides $|S| = p^m$, hence is a power of $p$.

\textit{There must be a fixed point.} Since $H$ is a Sylow $p$-subgroup, $p \nmid |G|/|H| = |G/H|$. Writing $|G/H|$ as a sum of orbit sizes, each of which is a power of $p$, the only way this sum is not divisible by $p$ is if at least one orbit has size $p^0 = 1$.

\textit{A fixed point gives conjugacy.} If $gH$ is a fixed point, then $s \cdot gH = gH$ for all $s \in S$, so $g^{-1}sg \in H$ for all $s \in S$. Hence $g^{-1}Sg \subseteq H$, and since $|g^{-1}Sg| = |S| = |H|$, we get $g^{-1}Sg = H$. So $S$ and $H$ are conjugate.
\end{proof}


\begin{corollary}[3.10]
A Sylow $p$-subgroup of $G$ is unique if and only if it is normal in $G$.
\end{corollary}

\begin{proof}
If $P$ is the unique Sylow $p$-subgroup, then $gPg^{-1} = P$ for all $g \in G$ (since conjugates are also Sylow $p$-subgroups), so $P \trianglelefteq G$. Conversely, if $P \trianglelefteq G$ then $P$ is its own conjugate, so by the Second Sylow Theorem it is the only Sylow $p$-subgroup.
\end{proof}

\begin{theorem}[3.9 — Third Sylow Theorem]
Let $G$ be a finite group and $p$ a prime. The number $n_p$ of Sylow $p$-subgroups satisfies $n_p \mid |G|$ and $n_p \equiv 1 \pmod{p}$.
\end{theorem}

\begin{proof}
\textit{$n_p \mid |G|$.} By the Second Sylow Theorem, all Sylow $p$-subgroups are conjugate to a fixed $K \in \mathrm{Syl}_p(G)$, so $n_p = |\mathrm{ccl}(K)| = |G : N_G(K)|$, which divides $|G|$.

\textit{$n_p \equiv 1 \pmod{p}$.} Let $P \in \mathrm{Syl}_p(G)$. Since conjugates of Sylow $p$-subgroups are again Sylow $p$-subgroups (by the Second Sylow Theorem), $P$ acts on the set $\mathrm{Syl}_p(G)$ by conjugation:
\[
  x \cdot Q := xQx^{-1}, \quad x \in P,\ Q \in \mathrm{Syl}_p(G).
\]
This action partitions $\mathrm{Syl}_p(G)$ into orbits. Let $F \subseteq \mathrm{Syl}_p(G)$ denote the set of fixed points (those $Q$ with $xQx^{-1} = Q$ for all $x \in P$), and let $\mathcal{O}$ denote the collection of orbits of size greater than 1. Then:
\[
  n_p = |\mathrm{Syl}_p(G)| = |F| + \sum_{O \in \mathcal{O}} |O|.
\]
By the Orbit-Stabiliser Theorem each $|O|$ divides $|P| = p^m$, so every $O \in \mathcal{O}$ has size $p^e$ for some $e \geq 1$, hence $p \mid |O|$. Therefore $n_p \equiv |F| \pmod{p}$.

It remains to show $F = \{P\}$, which gives $n_p \equiv 1 \pmod{p}$. Clearly $P \in F$. If $Q \in F$, then $xQx^{-1} = Q$ for all $x \in P$, so $P \leq N_G(Q)$. By Corollary 3.10, since $Q \trianglelefteq N_G(Q)$, $Q$ is the unique Sylow $p$-subgroup of $N_G(Q)$. But $P$ is also a Sylow $p$-subgroup of $N_G(Q)$, so $P = Q$. Hence $F = \{P\}$.
\end{proof}


\section*{Lecture 4 — Modules: the Basics}

\begin{definition}[4.1 — Left $R$-module]
Let $R$ be a ring. An abelian group $M$ is a \emph{left $R$-module} if there is a map $R \times M \to M$, $(r,m) \mapsto rm$, such that for all $r, r_1, r_2 \in R$ and $m, m_1, m_2 \in M$:
\begin{enumerate}[label=(\roman*)]
  \item $r(m_1 + m_2) = rm_1 + rm_2$;
  \item $(r_1 + r_2)m = r_1 m + r_2 m$;
  \item $(r_1 r_2)m = r_1(r_2 m)$;
  \item $1m = m$ if $1 \in R$.
\end{enumerate}
\end{definition}

\begin{lemma}[4.7]
Let $M$ be an $R$-module. Then for $a \in R$ and $m \in M$:
\begin{enumerate}[label=(\roman*)]
  \item $0m = 0$;
  \item $a0 = 0$;
  \item $(-a)m = -(am) = a(-m)$.
\end{enumerate}
\end{lemma}

\begin{proof}
(i) $0m = (0 + 0)m = 0m + 0m$. Subtracting $0m$ from both sides gives $0 = 0m$.

(ii) $a0 = a(0 + 0) = a0 + a0$. Subtracting $a0$ from both sides gives $0 = a0$.

(iii) By (i) and axiom (ii): $0 = 0m = (a + (-a))m = am + (-a)m$, so $(-a)m = -(am)$.
By (ii) and axiom (i): $0 = a0 = a(m + (-m)) = am + a(-m)$, so $a(-m) = -(am)$.
\end{proof}

\begin{definition}[4.8 — Submodule]
A subset $N$ of an $R$-module $M$ is an \emph{$R$-submodule} of $M$ if:
\begin{enumerate}[label=(\roman*)]
  \item $N \neq \emptyset$;
  \item $a - b \in N$ for all $a, b \in N$;
  \item $ra \in N$ for all $a \in N$, $r \in R$.
\end{enumerate}
\end{definition}

\begin{theorem}[4.12]
Let $(N_i)_{i \in I}$ be a family of $R$-submodules of an $R$-module $M$. Then $\bigcap_{i \in I} N_i$ is also an $R$-submodule of $M$.
\end{theorem}

\begin{proof}
Since $0 \in N_i$ for all $i$, we have $0 \in \bigcap N_i$, so it is non-empty. Let $x, y \in \bigcap N_i$ and $r \in R$. For each $i$, since $N_i$ is a submodule, $x - y \in N_i$ and $rx \in N_i$. Hence $x - y$ and $rx$ lie in $\bigcap N_i$.
\end{proof}

\begin{definition}[4.13 — Generated submodule, cyclic module]
Let $S \subseteq M$. The family $\mathcal{A} = \{N \mid N \text{ is an } R\text{-submodule of } M \text{ containing } S\}$ is non-empty (since $M \in \mathcal{A}$), so by Theorem 4.12 the intersection $\bigcap_{N \in \mathcal{A}} N$ is an $R$-submodule. It is the \emph{smallest} submodule of $M$ containing $S$, called the submodule \emph{generated by $S$} and denoted $(S)$. If $S = \{x_1, \ldots, x_n\}$ we write $(x_1, \ldots, x_n)$. An $R$-module $M$ is \emph{cyclic} if $M = (x)$ for some $x \in M$.
\end{definition}

\begin{theorem}[4.15]
Let $R$ be a ring with unity and $x_1, \ldots, x_n \in M$. Then
\[
  (x_1, \ldots, x_n) = \sum_{i=1}^n Rx_i := \{r_1x_1 + \cdots + r_nx_n \mid r_i \in R\}.
\]
\end{theorem}

\begin{proof}
$(\supseteq)$ The submodule $(x_1,\ldots,x_n)$ contains each $x_i$. By closure under scalar multiplication it contains $rx_i$ for all $r \in R$, and by closure under addition it contains every sum $r_1x_1 + \cdots + r_nx_n$. Hence $\sum Rx_i \subseteq (x_1,\ldots,x_n)$.

$(\subseteq)$ We show $\sum Rx_i$ is a submodule containing each $x_i$; then $(x_1,\ldots,x_n) \subseteq \sum Rx_i$ follows from minimality of $(x_1,\ldots,x_n)$.

Each $x_i \in \sum Rx_i$ since $x_i = 1 \cdot x_i$ (using unity). For the submodule conditions: $0 = \sum 0 \cdot x_i \in \sum Rx_i$, so it is non-empty. For $\sum r_ix_i, \sum s_ix_i \in \sum Rx_i$ and $r \in R$,
\[
  \sum r_ix_i - \sum s_ix_i = \sum(r_i - s_i)x_i \in \sum Rx_i, \qquad r\!\sum r_ix_i = \sum(rr_i)x_i \in \sum Rx_i. \qedhere
\]
\end{proof}


\begin{definition}[4.18 — Sum of submodules]
Let $N_1, \ldots, N_k$ be $R$-submodules of $M$. Their \emph{sum} $\sum_{i=1}^k N_i$ is the submodule of $M$ generated by $\bigcup_{i=1}^k N_i$, i.e.\ the smallest submodule of $M$ containing each $N_i$.
\end{definition}

\section*{Lecture 5 — Modules: continuation}

\begin{theorem}[4.19]

Let $N_1,\ldots,N_k$ be $R$-submodules of an $R$-module $M$. Then
\[
\sum_{i=1}^k N_i
=
\{x_1+\cdots+x_k \mid x_i\in N_i\}.
\]

\end{theorem}

\begin{proof}

Let
\[
A=\{x_1+\cdots+x_k\mid x_i\in N_i\}.
\]
It is straightforward to check that $A$ is a submodule of $M$, and it contains every $N_i$. Hence, by minimality,
\[
\sum_{i=1}^k N_i\subseteq A.
\]
Conversely, $\sum_{i=1}^k N_i$ contains every $N_i$ and is closed under addition, so it contains every element of $A$. Therefore
\[
A\subseteq \sum_{i=1}^k N_i.
\]

\end{proof}

\begin{definition}[4.21 — Direct sum]

The sum $N_1+\cdots+N_k$ is \emph{direct} if every element has a unique representation
\[
x=x_1+\cdots+x_k,\qquad x_i\in N_i.
\]
In this case we write
\[
N_1\oplus\cdots\oplus N_k.
\]

\end{definition}

\begin{theorem}[4.22]

Let $N_1,\ldots,N_k$ be submodules of an $R$-module $M$. The following are equivalent:

\begin{enumerate}[label=(\roman*)]

  \item $\sum_{i=1}^k N_i$ is direct;

  \item if $\sum_{i=1}^k x_i=0$, where $x_i\in N_i$, then every $x_i=0$;

  \item for every $i$,
  \[
  N_i\cap\sum_{j\neq i}N_j=\{0\}.
  \]

\end{enumerate}

\end{theorem}

\begin{proof}

(i) $\Rightarrow$ (ii): Zero already has the representation $0+\cdots+0$, so uniqueness implies that every $x_i=0$.

(ii) $\Rightarrow$ (iii): If
\[
x\in N_i\cap\sum_{j\neq i}N_j,
\]
then $x=\sum_{j\neq i}x_j$ for some $x_j\in N_j$. Hence
\[
x-\sum_{j\neq i}x_j=0,
\]
so (ii) gives $x=0$.

(iii) $\Rightarrow$ (i): Suppose
\[
\sum_{i=1}^k x_i=\sum_{i=1}^k y_i,
\qquad x_i,y_i\in N_i.
\]
For every $i$,
\[
x_i-y_i=-\sum_{j\neq i}(x_j-y_j)
\in N_i\cap\sum_{j\neq i}N_j=\{0\}.
\]
Thus $x_i=y_i$ for every $i$, so the representation is unique.

\end{proof}

\begin{definition}[4.23 — $R$-homomorphism]

Let $M$ and $N$ be $R$-modules. A map $f\colon M\to N$ is an \emph{$R$-homomorphism} if, for all $x,y\in M$ and $r\in R$,
\[
f(x+y)=f(x)+f(y),
\qquad
f(rx)=rf(x).
\]

The set of all $R$-homomorphisms $M\to N$ is denoted
\[
\operatorname{Hom}_R(M,N).
\]
An $R$-homomorphism $M\to M$ is called an \emph{endomorphism}, and
\[
\operatorname{End}_R(M):=\operatorname{Hom}_R(M,M).
\]

\end{definition}

\begin{theorem}[4.30]
Let $M = \bigoplus_{i=1}^k M_i$ be a direct sum of $R$-modules. Then
\[
  \mathrm{Hom}_R(M,M) \cong \bigl(\mathrm{Hom}_R(M_j,M_i)\bigr)_{1 \leq i,j \leq k}
\]
as rings, where the right-hand side is a ring $T$ of $k \times k$ matrices under the usual matrix addition and multiplication, with $(i,j)$ entry in $\mathrm{Hom}_R(M_j,M_i)$.
\end{theorem}

\begin{proof}
Let $\lambda_j : M_j \to M$ be the natural inclusion and $\pi_i : M \to M_i$ the natural projection. Two key identities hold:
\[
  \sum_{\ell=1}^k \lambda_\ell\pi_\ell = \mathrm{id}_M, \qquad \pi_p\lambda_q = \delta_{pq}\,\mathrm{id}_{M_p}.
\]
The first says every $x \in M$ equals the sum of its components; the second says projecting the $q$-th inclusion gives zero unless $p = q$.

Define $\sigma : \mathrm{Hom}_R(M,M) \to T$ by $\sigma(\phi) = (\pi_i\phi\lambda_j)_{ij}$.

\textit{Ring homomorphism.} Additivity is immediate. For multiplication, using the first identity:
\[
  (\sigma(\phi)\sigma(\psi))_{ij} = \sum_\ell \pi_i\phi\lambda_\ell\pi_\ell\psi\lambda_j = \pi_i\phi\!\left(\sum_\ell\lambda_\ell\pi_\ell\right)\!\psi\lambda_j = \pi_i\phi\,\mathrm{id}_M\,\psi\lambda_j = (\sigma(\phi\psi))_{ij}.
\]

\textit{Injective.} Suppose $\sigma(\phi) = 0$, so $\pi_i\phi\lambda_j = 0$ for all $i,j$. Applying $\lambda_i$ on the left and summing over $i$ using the first identity:
\[
  \phi\lambda_j = \mathrm{id}_M\cdot\phi\lambda_j = \Bigl(\sum_i\lambda_i\pi_i\Bigr)\phi\lambda_j = \sum_i \lambda_i(\pi_i\phi\lambda_j) = 0.
\]
Now applying $\pi_j$ on the right and summing over $j$:
\[
  \phi = \phi\cdot\mathrm{id}_M = \phi\Bigl(\sum_j\lambda_j\pi_j\Bigr) = \sum_j(\phi\lambda_j)\pi_j = 0.
\]

\textit{Surjective.} Given $f = (f_{ij}) \in T$, set $\phi = \sum_{i,j}\lambda_i f_{ij}\pi_j$. Using the second identity:
\[
  (\sigma(\phi))_{st} = \pi_s\phi\lambda_t = \sum_{i,j}\pi_s\lambda_i\,f_{ij}\,\pi_j\lambda_t = \sum_{i,j}\delta_{si}\,f_{ij}\,\delta_{jt} = f_{st}.
\]
Hence $\sigma(\phi) = f$, so $\sigma$ is surjective.
\end{proof}


\begin{definition}[Quotient module]
Let $N$ be an $R$-submodule of an $R$-module $M$. The set of cosets $M/N = \{a + N \mid a \in M\}$ becomes an $R$-module under
\[
  (a + N) + (b + N) = (a + b) + N, \qquad r(a + N) = ra + N.
\]
It is called the \emph{quotient module} of $M$ by $N$.
\end{definition}

\begin{theorem}[4.31 — Correspondence Theorem for Modules]
Let $N$ be an $R$-submodule of an $R$-module $M$. Every $R$-submodule of $M/N$ is of the form $U/N$, where $U$ is an $R$-submodule of $M$ containing $N$.
\end{theorem}

\begin{proof}
The canonical projection $f : M \to M/N$, $f(x) = x + N$, is an $R$-homomorphism:
\[
  f(x + y) = (x+y) + N = (x+N) + (y+N) = f(x) + f(y), \qquad f(rx) = rx + N = r(x+N) = rf(x).
\]

Let $X$ be an $R$-submodule of $M/N$. The structure of the proof is: define $U$ as the preimage of $X$ under $f$, show $U$ is a submodule of $M$ containing $N$, and conclude $X = U/N$.

Set $U = f^{-1}(X) = \{x \in M \mid f(x) \in X\}$.

\textit{$N \subseteq U$.} For any $x \in N$, $f(x) = x + N = 0 + N = 0_{M/N} \in X$ (since $X$ is a submodule, it contains $0$). Hence $x \in U$.

\textit{$U$ is an $R$-submodule.} Since $0 \in N \subseteq U$, $U$ is non-empty. For $x, y \in U$ and $r \in R$, since $f$ is an $R$-homomorphism:
\[
  f(x - y) = f(x) - f(y) \in X, \qquad f(rx) = rf(x) \in X,
\]
so $x - y \in U$ and $rx \in U$.

\textit{$X = U/N$.} By definition of $U$, $f(U) \subseteq X$. For the reverse inclusion, take any $x \in X$; since $f$ is surjective there exists $y \in M$ with $f(y) = x$, so $y \in U$ and $x \in f(U)$. Hence $X = f(U) = U/N$.
\end{proof}


\begin{theorem}[4.32 — First Isomorphism Theorem for Modules]
Let $f : M \to N$ be an $R$-homomorphism. Then $M/\ker f \cong f(M)$.
\end{theorem}

\begin{proof}
Define $g : M/\ker f \to f(M)$ by $g(x + \ker f) = f(x)$. We know from group theory that $g$ is an isomorphism of abelian groups (well-defined and injective because $x + \ker f = y + \ker f \iff x - y \in \ker f \iff f(x) = f(y)$; surjective onto $f(M)$ by definition). It remains to check that $g$ respects the $R$-action:
\[
  g(r(x + \ker f)) = g(rx + \ker f) = f(rx) = rf(x) = rg(x + \ker f). \qedhere
\]
\end{proof}

\begin{theorem}[Second Isomorphism Theorem for Modules]
Let $U$ and $V$ be $R$-submodules of an $R$-module $M$. Then $(U+V)/V \cong U/(U \cap V)$.
\end{theorem}

\begin{proof}
Define $f : U \to (U+V)/V$ by $f(u) = u + V$.

\textit{$f$ is an $R$-homomorphism.} For $u_1, u_2 \in U$ and $r \in R$:
\[
  f(u_1 + u_2) = (u_1+u_2)+V = (u_1+V)+(u_2+V) = f(u_1)+f(u_2),
\]
\[
  f(ru) = ru+V = r(u+V) = rf(u).
\]

\textit{$f$ is surjective.} Any element of $(U+V)/V$ has the form $(u+v)+V = u+V = f(u)$ for some $u \in U$, $v \in V$ (since $v \in V$ means $v + V = V$).

\textit{$\ker f = U \cap V$.} We have $f(u) = 0$ iff $u+V = V$ iff $u \in V$, so $\ker f = \{u \in U \mid u \in V\} = U \cap V$.

By the First Isomorphism Theorem, $U/(U \cap V) \cong (U+V)/V$.
\end{proof}


\begin{theorem}[Third Isomorphism Theorem for Modules]
Let $U \subseteq V$ be $R$-submodules of an $R$-module $M$. Then $(M/U)/(V/U) \cong M/V$.
\end{theorem}

\begin{proof}
Define $f : M/U \to M/V$ by $f(x+U) = x+V$.

\textit{$f$ is well-defined.} If $x+U = y+U$ then $x-y \in U \subseteq V$, so $x+V = y+V$.

\textit{$f$ is an $R$-homomorphism.} For $x+U, y+U \in M/U$ and $r \in R$:
\[
  f((x+U)+(y+U)) = f((x+y)+U) = (x+y)+V = (x+V)+(y+V) = f(x+U)+f(y+U),
\]
\[
  f(r(x+U)) = f(rx+U) = rx+V = r(x+V) = rf(x+U).
\]

\textit{$f$ is surjective.} For any $x+V \in M/V$, $f(x+U) = x+V$.

\textit{$\ker f = V/U$.} We have $f(x+U) = 0$ iff $x+V = V$ iff $x \in V$, so $\ker f = \{x+U \mid x \in V\} = V/U$.

By the First Isomorphism Theorem, $(M/U)/(V/U) \cong M/V$.
\end{proof}

\begin{theorem}[4.33]
Let $M_1, \ldots, M_n$ be $R$-modules and $A_i \leq M_i$ an $R$-submodule for each $i$. Then
\[
  \frac{M_1 \times \cdots \times M_n}{A_1 \times \cdots \times A_n} \cong \frac{M_1}{A_1} \times \cdots \times \frac{M_n}{A_n}.
\]
Since the product is finite, $\times$ and $\oplus$ agree here, so equivalently $(M_1 \oplus \cdots \oplus M_n)/(A_1 \oplus \cdots \oplus A_n) \cong M_1/A_1 \oplus \cdots \oplus M_n/A_n$.
\end{theorem}

\begin{proof}
Define $f : M_1 \times \cdots \times M_n \to M_1/A_1 \times \cdots \times M_n/A_n$ by
\[
  f(m_1, \ldots, m_n) = (m_1 + A_1, \ldots, m_n + A_n).
\]

\textit{$f$ is an $R$-homomorphism.} Each coordinate is the quotient map $\pi_i : M_i \to M_i/A_i$, and operations in a product are coordinatewise, so $f$ respects addition and scalar multiplication because each $\pi_i$ does.

\textit{$f$ is surjective.} Any element of the target is $(m_1 + A_1, \ldots, m_n + A_n)$ for some $m_i \in M_i$, and this equals $f(m_1, \ldots, m_n)$.

\textit{$\ker f = A_1 \times \cdots \times A_n$.} We have $f(m_1, \ldots, m_n) = 0$ iff $m_i + A_i = A_i$ for every $i$, iff $m_i \in A_i$ for every $i$.

By the First Isomorphism Theorem, $(M_1 \times \cdots \times M_n)/(A_1 \times \cdots \times A_n) \cong M_1/A_1 \times \cdots \times M_n/A_n$.
\end{proof}

\section*{Lecture 6 — Simple, Completely Reducible, and Free Modules}

\begin{definition}[5.1 — Simple module]
Let $R$ be a ring with unity. An $R$-module $M$ is called \emph{simple} (or \emph{irreducible}) if $M \neq 0$ and the only $R$-submodules of $M$ are $0$ and $M$.
\end{definition}

\begin{theorem}[5.5 — Characterisation of Simple Modules]
Let $R$ be a ring with unity and $M$ an $R$-module. The following are equivalent:
\begin{enumerate}[label=(\roman*)]
  \item $M$ is simple;
  \item $M \neq 0$, and $M$ is generated by any $0 \neq x \in M$;
  \item $M \cong R/I$, where $I$ is a maximal left ideal of $R$.
\end{enumerate}
\end{theorem}

\begin{proof}
(i) $\Rightarrow$ (ii): Let $0 \neq x \in M$. Then $Rx$ is a non-zero submodule of $M$, so $Rx = M$ by simplicity.

(ii) $\Rightarrow$ (iii): Take $0 \neq x \in M$ and define $f : R \to M$ by $f(r) = rx$. This is a surjective $R$-homomorphism, so letting $I = \ker f$, the First Isomorphism Theorem gives $R/I \cong M$. Note that (ii) implies $M$ is simple: if $0 \neq N \leq M$ and $0 \neq x \in N$, then $M = Rx \subseteq N$. Hence $R/I \cong M$ is simple. By the Correspondence Theorem (4.31), any left ideal $J$ with $I \subseteq J \subseteq R$ corresponds to a submodule $J/I$ of $R/I$. Simplicity forces $J/I = 0$ or $J/I = R/I$, i.e.\ $J = I$ or $J = R$, so $I$ is maximal.

(iii) $\Rightarrow$ (i): Let $0 \neq N \leq M \cong R/I$. By the Correspondence Theorem (4.31), $N \cong J/I$ for some left ideal $J$ with $I \subsetneq J \subseteq R$. Since $I$ is maximal, $J = R$, so $N = M$.
\end{proof}

\begin{theorem}[5.6 — Schur's Lemma]
Let $M$ be a simple $R$-module. Then $\mathrm{Hom}_R(M,M)$ is a division ring.
\end{theorem}

\begin{proof}
Let $0 \neq f \in \mathrm{Hom}_R(M,M)$. We show $f$ is invertible and $f^{-1} \in \mathrm{Hom}_R(M,M)$.

\textit{$f$ is bijective.} Since $\ker f$ and $\mathrm{im}\, f$ are submodules of $M$, simplicity forces each to be $0$ or $M$. As $f \neq 0$, we have $\ker f \neq M$ and $\mathrm{im}\, f \neq 0$, so $\ker f = 0$ (injective) and $\mathrm{im}\, f = M$ (surjective).

\textit{$f^{-1}$ is an $R$-homomorphism.} We cannot push $f^{-1}$ through operations directly, but we can push $f$ through them. The strategy is: to prove $f^{-1}(u) = v$, it suffices to show $f(u) = f(v)$, since $f$ is injective. For $x, y \in M$ and $r \in R$:
\begin{align*}
  f\!\left(f^{-1}(x+y)\right) &= x + y = f(f^{-1}(x)) + f(f^{-1}(y)) = f\!\left(f^{-1}(x) + f^{-1}(y)\right), \\
  f\!\left(f^{-1}(rx)\right)  &= rx = rf(f^{-1}(x)) = f\!\left(rf^{-1}(x)\right),
\end{align*}
where we used that $f$ is itself an $R$-homomorphism in the last step of each line. Injectivity of $f$ gives $f^{-1}(x+y) = f^{-1}(x) + f^{-1}(y)$ and $f^{-1}(rx) = rf^{-1}(x)$.
\end{proof}

\begin{definition}[5.7 — Completely reducible module]
Let $R$ be a ring with unity. An $R$-module $M$ is called \emph{completely reducible} (or \emph{semisimple}) if $M = \sum_{\alpha \in \Lambda} M_\alpha$, where each $M_\alpha$ is a simple $R$-submodule.
\end{definition}

\begin{theorem}[5.9]
Let $M = \sum_{\alpha \in \Lambda} M_\alpha$ be a sum of simple $R$-submodules, and let $K$ be a submodule of $M$. Then there exists $\Lambda' \subseteq \Lambda$ such that $M = K \oplus \bigl(\bigoplus_{\alpha \in \Lambda'} M_\alpha\bigr)$.
\end{theorem}

\begin{proof}
Let $S$ be the family of subsets $A \subseteq \Lambda$ satisfying two conditions: $\bigoplus_{\alpha \in A} M_\alpha$ is a direct sum, and $K \cap \bigoplus_{\alpha \in A} M_\alpha = 0$. Together these mean $K \oplus \bigoplus_{\alpha \in A} M_\alpha$ is a genuine direct sum. Since $\emptyset$ satisfies both conditions trivially, $S \neq \emptyset$. Partially order $S$ by inclusion; every chain $(A_i)$ in $S$ has upper bound $\bigcup A_i$ in $S$. By Zorn's Lemma, $S$ has a maximal element $A$.

Set $N = K \oplus \bigoplus_{\alpha \in A} M_\alpha$. We claim $N = M$. Fix any $\beta \in \Lambda$. Since $M_\beta$ is simple, $M_\beta \cap N$ is either $0$ or $M_\beta$. If $M_\beta \cap N = 0$, then $M_\beta$ intersects neither $K$ nor $\bigoplus_{\alpha \in A} M_\alpha$, so $A \cup \{\beta\}$ would satisfy both conditions in $S$, contradicting maximality of $A$. Hence $M_\beta \cap N = M_\beta$, i.e.\ $M_\beta \subseteq N$, for every $\beta \in \Lambda$. Since $M = \sum_\beta M_\beta$, we conclude $N = M$.
\end{proof}

\begin{corollary}[5.10]
If $M = \sum_{\alpha \in \Lambda} M_\alpha$ is a sum of simple $R$-submodules, then there exists $\Lambda' \subseteq \Lambda$ such that $M = \bigoplus_{\alpha \in \Lambda'} M_\alpha$.
\end{corollary}

\begin{proof}
Apply Theorem 5.9 with $K = \{0\}$.
\end{proof}

\begin{definition}[5.11 — Linear independence and basis]
A sequence $x_1, \ldots, x_n$ in an $R$-module $M$ is \emph{linearly independent} if $\sum_{i=1}^n a_i x_i = 0$ implies $a_1 = \cdots = a_n = 0$ for any $a_i \in R$. A subset $S \subseteq M$ is linearly independent if every finite sequence of distinct elements of $S$ is linearly independent.

A subset $B \subseteq M$ is a \emph{basis} of $M$ if $B$ generates $M$ and $B$ is linearly independent. An $R$-module that has a basis is called a \emph{free module}.
\end{definition}

\section*{Lecture 7 — Free Modules and Chain Conditions}

\begin{theorem}[5.19]
Let $M$ be a free $R$-module with basis $\{f_1, \ldots, f_n\}$. Then $M \cong R^n$ as $R$-modules.
\end{theorem}

\begin{proof}
Define $\varphi : M \to R^n$ by $\varphi\!\left(\sum_{i=1}^n r_i f_i\right) = (r_1, \ldots, r_n)$, i.e.\ read off the coordinates in the basis. Since every element of $M$ has a \emph{unique} representation as $\sum r_i f_i$ (by linear independence of the $f_i$), $\varphi$ is well-defined. It is an $R$-homomorphism: $\varphi\!\left(r\cdot\sum s_i f_i\right) = \varphi\!\left(\sum (rs_i)f_i\right) = (rs_1,\ldots,rs_n) = r\cdot(s_1,\ldots,s_n) = r\cdot\varphi\!\left(\sum s_i f_i\right)$, where the last scalar multiplication is coordinate-wise in $R^n$. It is injective (if $(r_1,\ldots,r_n) = 0$ then each $r_i = 0$, so $m = 0$), and surjective (every $(r_1,\ldots,r_n) \in R^n$ is the image of $\sum r_i f_i$).
\end{proof}

\begin{theorem}[5.20]
Let $M$ be a finitely generated free $R$-module. Then every basis of $M$ is finite.
\end{theorem}

\begin{proof}
Let $\{e_i\}_{i \in \Lambda}$ be a basis and $\{x_1, \ldots, x_n\}$ a finite generating set. Since elements of a module are finite sums by definition, each $x_j$ involves only finitely many nonzero coefficients, so is supported on a finite subset $S_j \subseteq \Lambda$. Every basis element must appear in the expression of some generator (otherwise it could not be reached), so $\Lambda \subseteq S_1 \cup \cdots \cup S_n$, a finite union of finite sets. Hence $\Lambda$ is finite.
\end{proof}

\begin{theorem}[5.21]
Let $M$ be a finitely generated free module over a commutative ring $R$. Then all bases of $M$ have the same number of elements.
\end{theorem}

\begin{proof}
Suppose $M$ has two bases of sizes $m$ and $n$ respectively. By Theorem 5.19, $M \cong R^m$ and $M \cong R^n$, so $R^m \cong R^n$. It suffices to show $m = n$.

Assume for contradiction $m < n$. Let $\phi : R^m \to R^n$ be an $R$-isomorphism with inverse $\psi = \phi^{-1}$. Let $(e_1,\ldots,e_m)$ and $(f_1,\ldots,f_n)$ be the standard bases of $R^m$ and $R^n$. Since $\phi$ is an $R$-homomorphism, it is determined by where it sends each $e_i$, and each image lands in $R^n$, so we can write
\[
  \phi(e_i) = \sum_{j=1}^n a_{ji} f_j, \qquad \psi(f_j) = \sum_{k=1}^m b_{kj} e_k.
\]
Let $A = (a_{ji})$ be the resulting $n \times m$ matrix and $B = (b_{kj})$ the $m \times n$ matrix. Since $\psi\phi = \mathrm{id}_{R^m}$ and $\phi\psi = \mathrm{id}_{R^n}$, substituting the expressions above and using linear independence of the bases gives $BA = I_m$ and $AB = I_n$.

Pad these to $n \times n$ matrices by setting $A' = \begin{pmatrix} A & 0 \end{pmatrix}$ and $B' = \begin{pmatrix} B \\ 0 \end{pmatrix}$. Then $A'B' = I_n$, so $\det(A'B') = 1$. But $B'A' = \begin{pmatrix} I_m & 0 \\ 0 & 0 \end{pmatrix}$, so $\det(B'A') = 0$. Since $R$ is commutative, determinants multiply, giving
\[
  1 = \det(A'B') = \det(A')\det(B') = \det(B'A') = 0,
\]
a contradiction. By symmetry $n < m$ also fails, so $m = n$.
\end{proof}

\begin{definition}[5.22 — Rank]
The number of elements in any basis of a finitely generated free module $M$ over a commutative ring $R$ is called the \emph{rank} of $M$, written $\mathrm{rk}(M)$. If $R$ is a field, the rank is called the \emph{dimension}.
\end{definition}

\begin{lemma}[5.24]
Every finitely generated $R$-module is a homomorphic image of a finitely generated free module.
\end{lemma}

\begin{proof}
Let $M$ be generated by $x_1, \ldots, x_n$. Define $\phi : R^n \to M$ by $\phi(r_1, \ldots, r_n) = \sum_{i=1}^n r_i x_i$. This is a surjective $R$-homomorphism (surjective since the $x_i$ generate $M$). The generators $x_i$ may satisfy relations, so $\phi$ need not be injective — but by the First Isomorphism Theorem $M \cong R^n / \ker\phi$, expressing $M$ as a quotient of the free module $R^n$.
\end{proof}

\begin{lemma}[5.25]
Let $R$ be a PID and $F$ a free $R$-module of rank $n$. Then every submodule $K \leq F$ is free of rank $m \leq n$.
\end{lemma}

\begin{proof}
By Theorem 5.19, $F \cong R^n$, so we work with $K \leq R^n$ and induct on $n$. The base case $n = 0$ is trivial since $K = 0$ is free on the empty set.

Define $\pi : R^n \to R$ by $\pi(x_1, \ldots, x_n) = x_1$, and restrict it to $K$.

\textit{Case 1: $\pi|_K = 0$.} Every element of $K$ has first coordinate zero, so $K \leq \{0\} \times R^{n-1} \cong R^{n-1}$. By induction $K$ is free of rank $\leq n-1 \leq n$.

\textit{Case 2: $\pi|_K \neq 0$.} The image $\pi(K)$ is a nonzero ideal of $R$. Since $R$ is a PID, $\pi(K) = Ra$ for some $a \neq 0$. Pick $k \in K$ with $\pi(k) = a$.

We claim $K = Rk \oplus \ker(\pi|_K)$, where $\ker(\pi|_K) = \{x \in K \mid x_1 = 0\}$ is the set of elements of $K$ with first coordinate zero, so $\ker(\pi|_K) \leq \{0\} \times R^{n-1} \cong R^{n-1}$.

\textit{$K = Rk + \ker(\pi|_K)$:} For any $x \in K$, write $\pi(x) = ba$ for some $b \in R$. Then $\pi(x - bk) = ba - ba = 0$, so $x - bk \in \ker(\pi|_K)$, and $x = bk + (x - bk)$.

\textit{The sum is direct:} If $ck \in \ker(\pi|_K)$, then $0 = \pi(ck) = ca$. Since $R$ is a PID (hence an integral domain) and $a \neq 0$, we get $c = 0$.

By induction $\ker(\pi|_K)$ is free of rank $s \leq n-1$, and $Rk \cong R$ is free of rank $1$. Hence $K \cong R^s \oplus R = R^{s+1}$, which has rank $m = s+1 \leq n$.
\end{proof}

\section*{Lecture 8 — Noetherian and Artinian Modules}

\begin{definition}[6.1 — Noetherian and Artinian modules]
An $R$-module $M$ is called \emph{Noetherian} if every ascending chain of submodules
\[
  M_1 \subseteq M_2 \subseteq M_3 \subseteq \cdots
\]
stabilises, i.e.\ there exists $k$ such that $M_k = M_{k+1} = \cdots$. This is called the \emph{ascending chain condition} (ACC). An $R$-module $M$ is called \emph{Artinian} if every descending chain of submodules stabilises (\emph{descending chain condition}, DCC).
\end{definition}

\begin{theorem}[6.5]
Let $M$ be an $R$-module. The following are equivalent:
\begin{enumerate}[label=(\roman*)]
  \item $M$ is Noetherian;
  \item every submodule of $M$ is finitely generated;
  \item every non-empty family $S$ of submodules of $M$ has a maximal element.
\end{enumerate}
\end{theorem}

\begin{proof}
(i) $\Rightarrow$ (ii): Let $N \leq M$ and suppose $N$ is not finitely generated. For any finite choice $a_1,\ldots,a_k \in N$ we have $N \neq (a_1,\ldots,a_k)$, so we can always pick $a_{k+1} \in N \setminus (a_1,\ldots,a_k)$. This produces an infinite strictly ascending chain $(a_1) \subsetneq (a_1,a_2) \subsetneq \cdots$, contradicting that $M$ is Noetherian.

(ii) $\Rightarrow$ (iii): Suppose $S$ has no maximal element. Then starting from any $N_0 \in S$ we can always find a strictly larger member, producing an infinite strictly ascending chain $N_0 \subsetneq N_1 \subsetneq N_2 \subsetneq \cdots$. Let $N = \bigcup_i N_i$. For $x, y \in N$ and $r \in R$, both $x$ and $y$ lie in some $N_\ell$ (take the larger index), so $x - y, rx \in N_\ell \subseteq N$; hence $N$ is a submodule. By (ii), $N = (a_1,\ldots,a_n)$ for some $a_j \in N$. Each $a_j$ lies in some $N_{u_j}$, so all generators lie in $N_\mu$ for $\mu = \max\{u_1,\ldots,u_n\}$. But then $N = (a_1,\ldots,a_n) \subseteq N_\mu \subseteq N$, so $N_\mu = N$, contradicting that the chain is strictly ascending.

(iii) $\Rightarrow$ (i): Given an ascending chain $M_1 \subseteq M_2 \subseteq \cdots$, by (iii) the family $\{M_i\}$ has a maximal element $M_k$. Then $M_k = M_{k+1} = \cdots$, so the chain stabilises.
\end{proof}

\begin{definition}[6.6 — Finitely cogenerated]
An $R$-module $M$ is \emph{finitely cogenerated} if for every family $(M_\alpha)_{\alpha \in \Lambda}$ of submodules of $M$,
\[
  \bigcap_{\alpha \in \Lambda} M_\alpha = 0 \implies \bigcap_{\alpha \in \Lambda'} M_\alpha = 0
\]
for some finite subset $\Lambda' \subseteq \Lambda$.
\end{definition}

\begin{theorem}[6.8]
Let $M$ be an $R$-module. The following are equivalent:
\begin{enumerate}[label=(\roman*)]
  \item $M$ is Artinian;
  \item every quotient module of $M$ is finitely cogenerated;
  \item every non-empty family $S$ of submodules of $M$ has a minimal element.
\end{enumerate}
\end{theorem}

\begin{proof}
(i) $\Rightarrow$ (ii): Let $N \leq M$ and take any family $(M_\alpha/N)_{\alpha \in \Lambda}$ of submodules of $M/N$ with $\bigcap_{\alpha \in \Lambda} M_\alpha/N = 0$, i.e.\ $\bigcap_{\alpha \in \Lambda} M_\alpha = N$. Consider the family of finite intersections $\bigl\{\bigcap_{\alpha \in \Lambda'} M_\alpha \mid \Lambda' \subseteq \Lambda \text{ finite}\bigr\}$. As we add more submodules to the intersection it can only get smaller, so this is a descending family of submodules of $M$. Since $M$ is Artinian every descending chain stabilises, so there exists a finite $\Lambda' \subseteq \Lambda$ such that
\[
  \bigcap_{\alpha \in \Lambda'} M_\alpha = \bigcap_{\alpha \in \Lambda} M_\alpha = N.
\]
Hence $\bigcap_{\alpha \in \Lambda'} M_\alpha/N = 0$, so $M/N$ is finitely cogenerated.

(ii) $\Rightarrow$ (iii): Suppose $S$ has no minimal element. Then starting from any $N_0 \in S$ we can always find a strictly smaller member, producing an infinite strictly descending chain $N_0 \supsetneq N_1 \supsetneq N_2 \supsetneq \cdots$. Let $N = \bigcap_i N_i$, which is a submodule of $M$. By (ii), $M/N$ is finitely cogenerated. The family $(N_i/N)$ of submodules of $M/N$ satisfies $\bigcap_i N_i/N = N/N = 0$, so by finite cogeneration there exists a finite subset, say $\{N_0,\ldots,N_k\}$, with $\bigcap_{i=0}^k N_i/N = 0$, i.e.\ $\bigcap_{i=0}^k N_i = N$. Since the chain is descending, $\bigcap_{i=0}^k N_i = N_k$, so $N_k = N \subseteq N_{k+1} \subsetneq N_k$, a contradiction.

(iii) $\Rightarrow$ (i): Given a descending chain $M_1 \supseteq M_2 \supseteq \cdots$, by (iii) the family $\{M_i\}$ has a minimal element $M_k$. Then $M_k = M_{k+1} = \cdots$, so the chain stabilises.
\end{proof}

\begin{theorem}[6.9]
Let $M$ be a Noetherian (respectively Artinian) $R$-module and $K \leq M$ a submodule. Then $K$ and $M/K$ are both Noetherian (respectively Artinian).
\end{theorem}

\begin{proof}
\textit{Noetherian case.} By Theorem 6.5, $M$ is Noetherian if and only if every submodule of $M$ is finitely generated. Every submodule of $K$ is also a submodule of $M$, hence finitely generated. By Theorem 6.5 again, $K$ is Noetherian. For $M/K$: by Theorem 4.31 (Correspondence Theorem), every submodule of $M/K$ is of the form $L/K$ for some submodule $L$ with $K \subseteq L \subseteq M$. Since $L$ is a submodule of $M$, it is finitely generated, and hence $L/K$ is finitely generated. By Theorem 6.5, $M/K$ is Noetherian.

\textit{Artinian case.} Any descending chain $K \supseteq L_1 \supseteq L_2 \supseteq \cdots$ of submodules of $K$ is also a descending chain in $M$, so it stabilises. Hence $K$ is Artinian. For $M/K$: a descending chain $M/K \supseteq L_1/K \supseteq L_2/K \supseteq \cdots$ of submodules of $M/K$ corresponds via Theorem 4.31 to a descending chain $M \supseteq L_1 \supseteq L_2 \supseteq \cdots$ in $M$, which stabilises. Hence $M/K$ is Artinian.
\end{proof}

\begin{theorem}[6.10]
Let $M$ be an $R$-module and $N \leq M$ a submodule. Then $M$ is Noetherian (respectively Artinian) if and only if both $N$ and $M/N$ are Noetherian (respectively Artinian).
\end{theorem}

\begin{proof}
$(\Rightarrow)$ This is Theorem 6.9.

$(\Leftarrow)$ \textit{Noetherian case.} Assume $N$ and $M/N$ are Noetherian. Let $K \leq M$ be any submodule; we show $K$ is finitely generated, which by Theorem 6.5 gives that $M$ is Noetherian.

By the Second Isomorphism Theorem, $K/(N \cap K) \cong (K+N)/N$. Since $(K+N)/N$ is a submodule of $M/N$ and $M/N$ is Noetherian, it is finitely generated by Theorem 6.5, and hence so is $K/(N \cap K)$. Say $K/(N \cap K)$ is generated by $\bar{x}_1, \ldots, \bar{x}_m$ with $x_i \in K$. Then
\[
  K = (x_1, \ldots, x_m) + (N \cap K).
\]
Since $N \cap K$ is a submodule of $N$ and $N$ is Noetherian, $N \cap K$ is finitely generated by Theorem 6.5, say by $y_1, \ldots, y_n \in N \cap K$. Therefore
\[
  K = (x_1, \ldots, x_m, y_1, \ldots, y_n),
\]
so $K$ is finitely generated and $M$ is Noetherian.

\textit{Artinian case.} Assume $N$ and $M/N$ are Artinian. Let $M \supseteq K_1 \supseteq K_2 \supseteq \cdots$ be a descending chain in $M$. Consider the two induced chains
\[
  K_1 \cap N \supseteq K_2 \cap N \supseteq \cdots \quad \text{in } N, \qquad \frac{K_1 + N}{N} \supseteq \frac{K_2 + N}{N} \supseteq \cdots \quad \text{in } M/N.
\]
Both stabilise (since $N$ and $M/N$ are Artinian), say $K_i \cap N = K_p \cap N$ for all $i \geq p$, and $(K_i+N)/N = (K_s+N)/N$ (i.e.\ $K_i + N = K_s + N$) for all $i \geq s$. Set $q = \max\{p, s\}$. Then for all $i \geq q$:
\[
  K_i \cap N = K_q \cap N \quad \text{and} \quad K_i + N = K_q + N.
\]
We claim $K_i = K_q$. Since the chain is descending and $i \geq q$, we already have $K_i \subseteq K_q$. For the reverse inclusion, take any $x \in K_q$. By the Second Isomorphism Theorem applied to $K_q$ and $N$:
\[
  K_q/(K_q \cap N) \cong (K_q + N)/N.
\]
In concrete terms, since $K_q + N = K_i + N$, every element of $K_q$ can be written as an element of $K_i$ plus an element of $N$: there exist $y \in K_i$ and $n \in N$ such that $x = y + n$. Then
\[
  n = x - y \in K_q,
\]
because $x \in K_q$ and $y \in K_i \subseteq K_q$. Hence $n \in K_q \cap N = K_i \cap N \subseteq K_i$. Therefore $x = y + n \in K_i$, giving $K_q \subseteq K_i$ and thus $K_i = K_q$. The original chain stabilises, so $M$ is Artinian.
\end{proof}


\begin{corollary}[6.11]
If $M_1, \ldots, M_n$ are Noetherian (respectively Artinian) $R$-modules, then $M_1 \oplus \cdots \oplus M_n$ is Noetherian (respectively Artinian).
\end{corollary}

\begin{proof}
By induction it suffices to treat the case of two modules $M$ and $N$. By Theorem 6.10, to show $M \oplus N$ is Noetherian (respectively Artinian) it suffices to find a single submodule of $M \oplus N$ such that both the submodule and its quotient module are Noetherian (respectively Artinian). Take $M \oplus 0 \leq M \oplus N$. Then:
\begin{itemize}
  \item the submodule $M \oplus 0 \cong M$ is Noetherian (respectively Artinian) by assumption;
  \item the quotient $(M \oplus N)/(M \oplus 0) \cong N$ is Noetherian (respectively Artinian) by assumption.
\end{itemize}
Hence by Theorem 6.10, $M \oplus N$ is Noetherian (respectively Artinian).
\end{proof}

\begin{definition}[6.12 — Noetherian and Artinian rings]
A ring $R$ is called \emph{left Noetherian} (respectively \emph{left Artinian}) if $R$, regarded as a left $R$-module, is Noetherian (respectively Artinian). By default, Noetherian (respectively Artinian) ring means left Noetherian (respectively left Artinian).
\end{definition}

\begin{theorem}[6.13]
Let $R$ be a ring. The following are equivalent:
\begin{enumerate}[label=(\roman*)]
  \item $R$ is Noetherian (respectively Artinian);
  \item every left ideal $A$ of $R$ is finitely generated (respectively every $R/A$ is finitely cogenerated);
  \item every non-empty set of left ideals of $R$ has a maximal (respectively minimal) element.
\end{enumerate}
\end{theorem}

\begin{proof}
This is Theorem 6.5 (respectively 6.8) applied to $R$ regarded as a left $R$-module, where submodules of $R$ are precisely the left ideals of $R$.
\end{proof}


\begin{theorem}[6.19]
Let $R_1, \ldots, R_n$ be Noetherian (respectively Artinian) rings. Then $R_1 \oplus \cdots \oplus R_n$ is Noetherian (respectively Artinian).
\end{theorem}

\begin{proof}
This follows immediately from Corollary 6.11, viewing each $R_i$ as an $R$-module.
\end{proof}

\medskip
\noindent\textbf{Product of ideals.} For ideals $A$ and $B$ in a ring $R$, define
\[
  AB = \left\{ \sum_{i=1}^n a_i b_i \;\middle|\; n \in \mathbb{N},\, a_i \in A,\, b_i \in B \right\}.
\]
Then $AB$ is again an ideal of $R$.

\begin{definition}[6.20 — Prime ideal]
Let $R$ be a ring. An ideal $P \neq R$ is \emph{prime} if whenever $A$ and $B$ are ideals in $R$ with $AB \subseteq P$, then $A \subseteq P$ or $B \subseteq P$.
\end{definition}

\noindent(For commutative $R$ this recovers the classical definition: $ab \in P \Rightarrow a \in P$ or $b \in P$.)

\begin{theorem}[6.21]
Let $R$ be a ring with unity. Then every maximal ideal is prime.
\end{theorem}

\begin{proof}
Let $M$ be a maximal ideal and let $A, B$ be ideals with $AB \subseteq M$. Suppose $A \not\subseteq M$. Since $M$ is maximal, the only ideal strictly containing $M$ is $R$ itself, so $A + M = R$. Write $1 = a + m$ for some $a \in A$ and $m \in M$. For any $b \in B$,
\[
  b = 1 \cdot b = (a + m)b = ab + mb \in AB + MB \subseteq M + M = M.
\]
Hence $B \subseteq M$, so $M$ is prime.
\end{proof}

\noindent\textit{Note:} The converse is false: $(0)$ is a prime ideal in $\mathbb{Z}$ but not a maximal ideal.

\section*{Lecture 9}

\begin{lemma}[6.22]
If $R$ is a Noetherian ring, then each ideal contains a finite product of prime ideals.
\end{lemma}

\begin{proof}
Let $\mathcal{F}$ be the family of ideals of $R$ that do not contain any finite product of prime ideals. We want to show $\mathcal{F} = \emptyset$.

Suppose for a contradiction that $\mathcal{F} \neq \emptyset$. Since $R$ is Noetherian, by Theorem 6.13(iii) the family $\mathcal{F}$ has a maximal element, say $A$.

We claim $A$ is not a prime ideal. Indeed, any prime ideal $P$ contains the one-factor product $P$ itself, so no prime ideal can belong to $\mathcal{F}$.

Since $A$ is not prime, there exist ideals $B$ and $C$ with $BC \subseteq A$, but $B \not\subseteq A$ and $C \not\subseteq A$. Then $B + A \supsetneq A$ and $C + A \supsetneq A$, so by maximality of $A$ in $\mathcal{F}$, neither $B + A$ nor $C + A$ belongs to $\mathcal{F}$. Hence both contain finite products of prime ideals. But then $(B+A)(C+A)$ also contains a finite product of prime ideals, and
\[
  (B+A)(C+A) \subseteq BC + BA + AC + A^2 \subseteq A.
\]
So $A$ contains a finite product of prime ideals, contradicting $A \in \mathcal{F}$. Hence $\mathcal{F} = \emptyset$.
\end{proof}

\begin{theorem}[6.23 — Hilbert's Basis Theorem]
Let $R$ be a Noetherian ring. Then the polynomial ring $R[x]$ is also Noetherian.
\end{theorem}

\begin{proof}
Suppose for a contradiction that $I \subseteq R[x]$ is a left ideal that is not finitely generated. We recursively construct a sequence of polynomials as follows. Set
\[
  J_0 = (0) \subseteq R[x].
\]
Since $I$ is not finitely generated, $I \setminus J_0 \neq \emptyset$. Pick $f_0 \in I \setminus J_0$ of minimal degree (exists since degrees are non-negative integers). For $n \geq 1$, given $f_0, \ldots, f_{n-1}$, let
\[
  J_n = (f_0, \ldots, f_{n-1}) \subseteq R[x]
\]
be the left ideal in $R[x]$ generated by $f_0, \ldots, f_{n-1}$. Since $I$ is not finitely generated, $I \setminus J_n \neq \emptyset$, so we may pick $f_n \in I \setminus J_n$ of minimal degree.

Since $J_n \subsetneq J_{n+1}$, the set $I \setminus J_{n+1} \subsetneq I \setminus J_n$, so $\deg(f_n)$ is non-decreasing in $n$.

For each $n$, let $a_n \in R$ denote the leading coefficient of $f_n$. Consider the ascending chain of left ideals in $R$:
\[
  L_1 \subseteq L_2 \subseteq \cdots, \qquad \text{where } L_n = (a_0, \ldots, a_{n-1}) \subseteq R.
\]
Since $R$ is Noetherian, this chain stabilises: there exists $N \in \mathbb{N}$ such that $L_N = L_{N+1}$, that is,
\[
  a_N \in L_N = (a_0, \ldots, a_{N-1}) \subseteq R.
\]
So we may write $a_N = \sum_{i < N} u_i a_i$ for some $u_i \in R$. Now define
\[
  g = \sum_{i < N} u_i x^{\deg(f_N) - \deg(f_i)} f_i \in J_N.
\]
The exponents $\deg(f_N) - \deg(f_i) \geq 0$ since degrees are non-decreasing. The leading term of $u_i x^{\deg(f_N)-\deg(f_i)} f_i$ is $u_i a_i x^{\deg(f_N)}$, so the leading term of $g$ is
\[
  \left(\sum_{i<N} u_i a_i\right) x^{\deg(f_N)} = a_N x^{\deg(f_N)},
\]
which equals the leading term of $f_N$. Hence $f_N - g$ has degree strictly less than $\deg(f_N)$.

Since $f_N \notin J_N$ and $g \in J_N$, we have $f_N - g \notin J_N$. Since $f_N, g \in I$, we have $f_N - g \in I$. Thus $f_N - g \in I \setminus J_N$.

\begin{itemize}
  \item If $\deg(f_N) > 0$: then $f_N - g \in I \setminus J_N$ has degree strictly less than $\deg(f_N)$, contradicting the minimality of $\deg(f_N)$.
  \item If $\deg(f_N) = 0$: then all degrees are $0$, so $f_N - g = 0$, meaning $f_N = g \in J_N$, contradicting $f_N \in I \setminus J_N$.
\end{itemize}

In both cases we reach a contradiction, so $I$ must be finitely generated and $R[x]$ is Noetherian.
\end{proof}

\begin{definition}[7.1 — Nilpotent elements and ideals]
Let $R$ be a ring.
\begin{itemize}
  \item An element $a \in R$ is \emph{nilpotent} if $a^n = 0$ for some $n \in \mathbb{N}$.
  \item A left ideal $I$ of $R$ is a \emph{nil ideal} if every element of $I$ is nilpotent.
  \item A left ideal $I$ of $R$ is \emph{nilpotent} if $I^n = 0$ for some $n \in \mathbb{N}$, where $I^n$ denotes the set of all finite sums of products $i_1 i_2 \cdots i_n$ with each $i_k \in I$.
\end{itemize}
Every nilpotent ideal is nil (if $I^n = 0$ then any $a \in I$ satisfies $a^n \in I^n = 0$), but the converse fails in general.
\end{definition}

\begin{theorem}[7.5]
If $J$ is a nil left ideal in an Artinian ring $R$, then $J$ is nilpotent.
\end{theorem}

\begin{proof}
Suppose for a contradiction that $J^k \neq 0$ for all $k \in \mathbb{N}$. The family $(J, J^2, J^3, \ldots)$ is a descending chain of left ideals of $R$. Since $R$ is Artinian, this family has a minimal element, say $B = J^m$ for some $m$. Then
\[
  B^2 = J^{2m} \subseteq J^m = B,
\]
and by minimality of $B$ we have $B^2 = B$.

Now consider the family
\[
  \mathcal{F} = \{A \mid A \text{ is a left ideal of } R,\; A \subseteq B,\; BA \neq 0\}.
\]
Since $B^2 = B \neq 0$, we have $B \in \mathcal{F}$, so $\mathcal{F} \neq \emptyset$. Since $R$ is Artinian, $\mathcal{F}$ has a minimal element $A$.

Since $BA \neq 0$, there exists $a \in A$ with $Ba \neq 0$. Note that $Ba = \{ba \mid b \in B\}$ is a left ideal contained in $A$, and
\[
  B(Ba) = B^2 a = Ba \neq 0,
\]
so $Ba \in \mathcal{F}$. By minimality of $A$ we conclude $Ba = A$.

Hence there exists $b \in B$ with $ba = a$. Applying $b$ repeatedly gives $b^i a = a$ for all $i \in \mathbb{N}$. But $b \in B \subseteq J$ and $J$ is nil, so $b^n = 0$ for some $n \in \mathbb{N}$. Then $a = b^n a = 0$, contradicting $Ba \neq 0$.

Therefore $J^k = 0$ for some $k \in \mathbb{N}$, i.e.\ $J$ is nilpotent.
\end{proof}

\begin{lemma}[7.6]
Let $R$ be a Noetherian ring. Then the sum of nilpotent left ideals in $R$ is a nilpotent left ideal.
\end{lemma}

\begin{proof}
Let $B = \sum_{i \in \Lambda} A_i$ be a sum of nilpotent left ideals. Since $R$ is Noetherian, $B$ is finitely generated as a left ideal, say $B = (x_1, \ldots, x_m)$. Each $x_j$ lies in a finite sum of the $A_i$'s, so $B$ is contained in a finite sum $A_1 + \cdots + A_n$ (reindexing if necessary), and hence $B = A_1 + \cdots + A_n$.

It remains to show $B$ is nilpotent. For each $i$, let $n_i \in \mathbb{N}$ be such that $A_i^{n_i} = 0$. Set $N = n_1 + \cdots + n_n$. Any product of $N$ elements from $B = A_1 + \cdots + A_n$, when expanded, is a sum of products of the form $a_{i_1} a_{i_2} \cdots a_{i_N}$ with each $a_{i_k} \in A_{i_k}$. By pigeonhole, among the $N$ indices $i_1, \ldots, i_N \in \{1, \ldots, n\}$, at least one value $j$ appears at least $n_j$ times. The consecutive factors from $A_j$ multiply to give an element of $A_j^{n_j} = 0$. Hence every such product is $0$, so $B^N = 0$.
\end{proof}

\section*{Lecture 10}

\begin{theorem}[7.7]
Let $R$ be a Noetherian ring having no non-zero nilpotent ideals. Then $R$ has no non-zero nil ideals.
\end{theorem}

\noindent\textit{(Proof non-examinable.)}


\begin{theorem}[7.8]
Let $N$ be a nil ideal in a Noetherian ring $R$. Then $N$ is nilpotent.
\end{theorem}

\begin{proof}
Let $T$ be the sum of all nilpotent ideals in $R$. By Lemma 7.6, $T$ is itself nilpotent, say $T^k = 0$.

We claim $R/T$ has no non-zero nilpotent ideals. Suppose $A/T$ is a nilpotent ideal in $R/T$, say $(A/T)^m = 0$. Then $(A^m + T)/T = 0$, so $A^m \subseteq T$. Since $T^k = 0$, we get $A^{mk} \subseteq T^k = 0$, so $A$ is nilpotent. Hence $A \subseteq T$, giving $A/T = 0$.

Since $R/T$ is Noetherian (a quotient of a Noetherian ring) and has no non-zero nilpotent ideals, Theorem 7.7 gives that $R/T$ has no non-zero nil ideals. Consider $(N+T)/T$: every element has the form $n + T$ for some $n \in N$ (since $t + T = T$ for $t \in T$), and since $N$ is nil there exists $k$ with $n^k = 0$, giving $(n+T)^k = n^k + T = T$. Hence $(N+T)/T$ is a nil ideal in $R/T$, so $(N+T)/T = 0$, which means $N \subseteq T$. Since $T$ is nilpotent, $N$ is nilpotent.
\end{proof}

\begin{lemma}[7.9]
Let $A$ be a minimal left ideal in a ring $R$. Then either $A^2 = 0$ or $A = Re$ for some idempotent $e \in R$.
\end{lemma}

\begin{proof}
Assume $A^2 \neq 0$.

\textit{Step 1: Find $e \in A$ with $ea = a$.}
Since $A^2 \neq 0$, there exists $a \in A$ with $Aa \neq 0$. The set $Aa = \{ca \mid c \in A\}$ is a left ideal contained in $A$ (since $A$ is a left ideal and $a \in A$). By minimality of $A$, we have $Aa = A$. In particular, there exists $e \in A$ with $ea = a$.

\textit{Step 2: Show $e^2 = e$.}
We compute $(e^2 - e)a = e(ea) - ea = ea - ea = 0$. To conclude $e^2 = e$, define the left annihilator of $a$ in $A$:
\[
  \mathrm{Ann}_A(a) := \{c \in A \mid ca = 0\}.
\]
This is a left ideal of $R$ contained in $A$. Since $ea = a \neq 0$, we have $e \notin \mathrm{Ann}_A(a)$, so $\mathrm{Ann}_A(a) \neq A$. By minimality of $A$, we conclude $\mathrm{Ann}_A(a) = 0$. Since $e^2 - e \in A$ and $(e^2 - e)a = 0$, we get $e^2 - e \in \mathrm{Ann}_A(a) = 0$, hence $e^2 = e$.

\textit{Step 3: Show $A = Re$.}
Since $e \in A$ and $A$ is a left ideal, $Re \subseteq A$. Since $e \neq 0$ (as $ea = a \neq 0$), we have $Re \neq 0$. By minimality of $A$, $Re = A$.
\end{proof}

\begin{theorem}[7.10 — Wedderburn--Artin]
Let $R$ be a left (or right) Artinian ring with unity and no non-zero nilpotent ideals. Then
\[
  R \cong M_{n_1}(D_1) \oplus \cdots \oplus M_{n_k}(D_k)
\]
for some $n_1, \ldots, n_k \in \mathbb{N}$ and division rings $D_1, \ldots, D_k$.
\end{theorem}

\begin{proof}
We build the result from two claims, then combine them.

\medskip
\noindent\textbf{Claim 1.} \textit{Each non-zero left ideal $A$ of $R$ is of the form $Re$ for some idempotent $e \in R$.}

\smallskip
By Lemma~7.9, every minimal left ideal is either nilpotent or of the form $Re$ for some idempotent $e$. Since $R$ has no non-zero nilpotent ideals, every minimal left ideal is of the form $Re$.

Now let $A$ be any non-zero left ideal. Since $R$ is Artinian, $A$ contains a minimal left ideal, hence a non-zero idempotent. Thus the family
\[
  \mathcal{F} = \{ R(1-f) \cap A \mid f \text{ a non-zero idempotent in } A \}
\]
is non-empty, and by the Artinian condition has a minimal member $R(1-e) \cap A$.

\medskip
\noindent\textbf{Subclaim} \textit{(non-examinable).} $R(1-e) \cap A = 0$.

\medskip
Invoking the subclaim: for any $a \in A$, $a(1-e) \in R(1-e) \cap A = 0$, so $a = ae \in Re$. Hence $A \subseteq Re \subseteq A$, giving $A = Re$.

\medskip
\noindent\textbf{Claim 2.} \textit{$R = Re_1 \oplus \cdots \oplus Re_n$ is a finite direct sum of minimal left ideals.}

\smallskip
Let $S$ be the sum of all minimal left ideals in $R$; since $R$ is Artinian it has at least one minimal left ideal, so $S \neq 0$. By Claim~1, $S = Re$ for some idempotent $e$.

Suppose $R(1-e) \neq 0$. Then $R(1-e)$ is a non-zero left ideal, so by the Artinian condition it contains a minimal left ideal $A$. But $A \subseteq S = Re$, giving $A \subseteq Re \cap R(1-e)$. For any $x \in Re \cap R(1-e)$, write $x = ae = b(1-e)$ for some $a, b \in R$. Then $xe = ae^2 = ae = x$, and also $xe = b(1-e)e = b(e - e^2) = 0$, so $x = 0$. Hence $Re \cap R(1-e) = 0$, giving $A = 0$, a contradiction.

Therefore $R(1-e) = 0$. Since $R$ has unity, $1 - e = 0$, i.e.\ $e = 1$, so $S = R$.

Thus $R$ is the sum of all its minimal left ideals. By Theorem~5.9 there exists a subfamily $(A_i)_{i \in \Lambda'}$ such that $R = \bigoplus_{i \in \Lambda'} A_i$. To see this is finite, write $1 = e_{i_1} + \cdots + e_{i_n}$ as a finite sum with $e_{i_j} \in A_{i_j}$. Then for any $r \in R$, $r = r \cdot 1 = re_{i_1} + \cdots + re_{i_n}$, so $R = Re_{i_1} \oplus \cdots \oplus Re_{i_n}$. After reindexing, $R = Re_1 \oplus \cdots \oplus Re_n$.

\medskip
\noindent\textbf{Conclusion.} The plan: compute $\mathrm{End}_R(R) = \mathrm{Hom}_R(R,R)$ as a matrix ring using Theorem~4.30, use Schur's Lemma to see that it is block diagonal, and finally recover $R$ from $\mathrm{End}_R(R)$.

\smallskip
\textit{Step 1: Group isomorphic summands.}
By Claim~2, $R = Re_1 \oplus \cdots \oplus Re_n$ as left $R$-modules, where each $Re_i$ is a minimal left ideal, i.e.\ a simple left $R$-module. Sort the summands into isomorphism classes; say there are $k$ classes, of sizes $n_1, \ldots, n_k$, and choose a representative $N_j$ of the $j$-th class. Thus $N_i \not\cong N_j$ for $i \neq j$, and every $Re_\ell$ is isomorphic to exactly one $N_j$. The order of summands in a direct sum is irrelevant, so we may reorder them so that each class is adjacent. Replacing each summand by its (isomorphic) representative then gives an isomorphism of left $R$-modules
\[
  R \;\cong\; M := \underbrace{(N_1 \oplus \cdots \oplus N_1)}_{n_1} \oplus \cdots \oplus \underbrace{(N_k \oplus \cdots \oplus N_k)}_{n_k}.
\]
A module isomorphism $\theta : R \to M$ induces a ring isomorphism $\mathrm{End}_R(R) \to \mathrm{End}_R(M)$, $\varphi \mapsto \theta\varphi\theta^{-1}$. Hence $\mathrm{End}_R(R) \cong \mathrm{End}_R(M)$, and it suffices to compute $\mathrm{End}_R(M)$.

\smallskip
\textit{Step 2: Apply Schur.}
Since each $N_j$ is simple, Schur's Lemma~5.6 says every non-zero homomorphism between $N_i$ and $N_j$ is an isomorphism. Therefore
\[
  \mathrm{Hom}_R(N_i, N_j) = 0 \quad (i \neq j), \qquad D_j := \mathrm{End}_R(N_j) \text{ is a division ring.}
\]

\smallskip
\textit{Step 3: Compute $\mathrm{End}_R(M)$.}
View $M$ as a direct sum of $n = n_1 + \cdots + n_k$ simple summands. By Theorem~4.30, $\mathrm{End}_R(M)$ is isomorphic to the ring of $n \times n$ matrices whose $(p,q)$ entry lies in $\mathrm{Hom}_R(\text{$q$-th summand}, \text{$p$-th summand})$. By Step~2, this entry is $0$ if the two summands are copies of different $N_i, N_j$, and is an arbitrary element of $D_j$ if both are copies of $N_j$. Since each class is adjacent, these are exactly the block-diagonal matrices
\[
  \begin{pmatrix}
    A_1 & & 0 \\
    & \ddots & \\
    0 & & A_k
  \end{pmatrix}, \qquad A_j \in M_{n_j}(D_j).
\]
Block-diagonal matrices add and multiply block by block, $\mathrm{diag}(A_1,\ldots,A_k)\,\mathrm{diag}(B_1,\ldots,B_k) = \mathrm{diag}(A_1B_1,\ldots,A_kB_k)$, so sending a matrix to its tuple of blocks is a ring isomorphism. Hence
\[
  \mathrm{End}_R(R) \;\cong\; \mathrm{End}_R(M) \;\cong\; M_{n_1}(D_1) \oplus \cdots \oplus M_{n_k}(D_k).
\]

\smallskip
\textit{Step 4: Identify $\mathrm{End}_R(R)$ with $R^{\mathrm{op}}$.}
Every $\varphi \in \mathrm{End}_R(R)$ is determined by $b := \varphi(1)$, since $\varphi(r) = \varphi(r \cdot 1) = r\varphi(1) = rb$. So every endomorphism is right multiplication $\rho_b : r \mapsto rb$ for a unique $b \in R$, and conversely each $\rho_b$ is an endomorphism. Composition reverses the order of multiplication: $(\rho_b \circ \rho_c)(r) = rcb$, so $\rho_b \circ \rho_c = \rho_{cb} = \rho_{b \cdot_{\mathrm{op}} c}$. Hence $b \mapsto \rho_b$ is a ring isomorphism $R^{\mathrm{op}} \cong \mathrm{End}_R(R)$, and combining with Step~3,
\[
  R^{\mathrm{op}} \;\cong\; M_{n_1}(D_1) \oplus \cdots \oplus M_{n_k}(D_k).
\]

\smallskip
\textit{Step 5: Pass to opposite rings.}
Taking opposites of both sides (using $(R^{\mathrm{op}})^{\mathrm{op}} = R$), and noting that the opposite of a direct sum is the direct sum of the opposites,
\[
  R \;\cong\; M_{n_1}(D_1)^{\mathrm{op}} \oplus \cdots \oplus M_{n_k}(D_k)^{\mathrm{op}}.
\]
Transposition gives $M_n(D)^{\mathrm{op}} \cong M_n(D^{\mathrm{op}})$, and $D^{\mathrm{op}}$ is again a division ring. Relabelling $D_j^{\mathrm{op}}$ as $D_j$, we conclude
\[
  R \;\cong\; M_{n_1}(D_1) \oplus \cdots \oplus M_{n_k}(D_k). \qedhere
\]
\end{proof}

\begin{theorem}[7.11]
Let $R$ be a left (or right) Artinian ring with unity and no non-zero nilpotent ideals. Then $R$ is also right and left Artinian, and right and left Noetherian.
\end{theorem}

\noindent\textit{(Non-examinable. Follows from Wedderburn--Artin and the fact that matrix rings over division rings are Noetherian and Artinian on both sides.)}

\begin{theorem}[7.13 — Maschke]
If $G$ is a finite group, then
\[
  \mathbb{C}[G] \cong M_{n_1}(\mathbb{C}) \oplus \cdots \oplus M_{n_k}(\mathbb{C})
\]
for some $n_1, \ldots, n_k \in \mathbb{N}$.
\end{theorem}

\noindent\textit{(Non-examinable. Follows from Wedderburn--Artin: one shows $\mathbb{C}[G]$ has no non-zero nilpotent ideals and is Artinian, then uses that the only division ring containing $\mathbb{C}$ in its centre as an algebraically closed field is $\mathbb{C}$ itself.)}

\begin{theorem}[7.14]
Let $R$ be a ring with unity. The following are equivalent:
\begin{enumerate}[label=(\roman*)]
  \item $R$ is left Artinian with no non-zero nilpotent ideals;
  \item $R$ is left Artinian with no non-zero nil ideals;
  \item $R$ is a finite direct sum of minimal left ideals;
  \item every left ideal of $R$ is of the form $Re$ for some idempotent $e$;
  \item $R$ is a finite direct sum of matrix rings over division rings.
\end{enumerate}
\end{theorem}

\noindent\textit{(Non-examinable.)}

\begin{definition}[7.15 — Semisimple Artinian ring]
A ring satisfying any of the equivalent conditions of Theorem 7.14 is called a \emph{semisimple Artinian ring}.
\end{definition}

\section*{Lecture 11 — Smith Normal Form over a PID}

Throughout this chapter $R$ is a principal ideal domain with unity.

\begin{definition}[8.1 — Row module, column module]
Let $A \in M_{m \times n}(R)$. The submodule of $R^n$ generated by the $m$ rows of $A$ is called the \emph{row module} of $A$, denoted $\mathcal{R}(A)$. The submodule of $R^m$ generated by the $n$ columns of $A$ is called the \emph{column module} of $A$, denoted $\mathcal{C}(A)$.

Since $\mathcal{R}(A)$ and $\mathcal{C}(A)$ are finitely generated submodules of the free modules $R^n$ and $R^m$ respectively, both are free by Lemma~5.25.
\end{definition}

\begin{definition}[8.2 — Row rank, column rank]
Let $A \in M_{m \times n}(R)$. The rank of $\mathcal{R}(A)$ is called the \emph{row rank} of $A$, and the rank of $\mathcal{C}(A)$ is called the \emph{column rank} of $A$.
\end{definition}

\begin{theorem}[8.3]
Let $A \in M_{m \times n}(R)$ and let $P$, $Q$ be invertible $m \times m$ and $n \times n$ matrices over $R$ respectively. Then row rank and column rank of $PAQ$ equal those of $A$.
\end{theorem}

\noindent\textit{(Proof non-examinable.)}

\begin{definition}[8.4/8.6 — Elementary operations and matrices]
Let $A \in M_{m \times n}(R)$. There are three \emph{elementary row operations}, each achieved by left-multiplying $A$ by an invertible \emph{elementary matrix} (column operations are achieved by right-multiplying):
\begin{enumerate}[label=(\roman*)]
  \item \textit{Swap} rows $i$ and $j$: multiply on the left by $S_{ij} = I - E_{ii} - E_{jj} + E_{ij} + E_{ji}$. Here $S_{ij}^{-1} = S_{ij}$.
  \item \textit{Scale} row $i$ by invertible $\alpha \in R$: multiply on the left by $M_i(\alpha) = I + (\alpha-1)E_{ii}$. Here $M_i(\alpha)^{-1} = M_i(\alpha^{-1})$.
  \item \textit{Add} $\alpha \in R$ times row $j$ to row $i$: multiply on the left by $C_{ij}(\alpha) = I + \alpha E_{ij}$. Here $C_{ij}(\alpha)^{-1} = C_{ij}(-\alpha)$.
\end{enumerate}
Since each elementary matrix is invertible, any sequence of elementary row and column operations transforms $A$ into an equivalent matrix $PAQ$.
\end{definition}

\begin{definition}[8.7 — Matrix equivalence]
Two matrices $A, B \in M_{m \times n}(R)$ are \emph{equivalent} if there exist invertible matrices $P \in M_m(R)$ and $Q \in M_n(R)$ such that $B = PAQ$.
\end{definition}

\begin{theorem}[8.8 — Smith Normal Form]
Let $A \in M_{m \times n}(R)$ where $R$ is a PID. Then $A$ is equivalent to a matrix of the form
\[
  \begin{pmatrix} a_1 & & & \\ & \ddots & & \\ & & a_r & \\ & & & 0 \end{pmatrix},
\]
where $a_1, \ldots, a_r \neq 0$ and $a_1 \mid a_2 \mid \cdots \mid a_r$.
\end{theorem}

\begin{proof}
Define the \emph{length} $\ell(a)$ of a non-zero $a \in R$ to be the number of prime factors in the factorisation $a = p_1 \cdots p_k$ counted with multiplicity, with the convention $\ell(u) = 0$ for units.

\textit{Step 1: bring a minimal-length entry to position $(1,1)$.}
If $A = 0$ the result is immediate. Otherwise, among all non-zero entries pick $a_{ij}$ with minimal $\ell(a_{ij})$. Row and column swaps bring it to position $(1,1)$.

\textit{Step 2: clear the first row.}
Fix any column $k \geq 2$. If $a_{11} \mid a_{1k}$, write $a_{1k} = a_{11}q$ and apply $C_k - qC_1$, zeroing the $(1,k)$ entry without changing $a_{11}$. If $a_{11} \nmid a_{1k}$, let $d = \gcd(a_{11}, a_{1k})$ and write $a_{11} = ds$, $a_{1k} = dt$. By Bézout there exist $u, v \in R$ with $us + vt = 1$, so the matrix $\bigl(\begin{smallmatrix} u & t \\ v & -s \end{smallmatrix}\bigr)$ is invertible (determinant $-1$). Right-multiplying $A$ by the $n \times n$ matrix with this $2 \times 2$ block in the $(1,k)$-positions and the identity elsewhere replaces the $(1,1)$ entry by $a_{11}u + a_{1k}v = d$ and the $(1,k)$ entry by $a_{11}t - a_{1k}s = 0$. Since $d$ properly divides $a_{11}$, we have $\ell(d) < \ell(a_{11})$. Repeating across all columns clears the first row.

\textit{Step 3: clear the first column.}
The identical argument applied to rows (left-multiplying by the analogous block matrices) clears all entries below $(1,1)$.

\textit{Step 4: termination.}
Clearing the first column may reintroduce non-zero entries in the first row, requiring another pass of Step 2. However this only occurs when $a_{11}$ does not divide some first-column entry, which triggers the Bézout case and strictly reduces $\ell(a_{11})$. Since $\ell(a_{11}) \geq 0$ is a non-negative integer, the alternating process terminates after at most $\ell(a_{11})$ restarts, yielding
\[
  P_1 A Q_1 = \begin{pmatrix} a_1 & 0 & \cdots & 0 \\ 0 & & & \\ \vdots & & A_1 & \\ 0 & & & \end{pmatrix},
\]
where $A_1$ is an $(m-1)\times(n-1)$ submatrix.

\textit{Step 5: induction.}
Apply the same procedure to $A_1$ and continue. By induction on $m + n$ we obtain $PAQ = \mathrm{diag}(a_1, \ldots, a_r, 0, \ldots, 0)$.

\textit{Step 6: enforce $a_1 \mid a_2 \mid \cdots \mid a_r$.}
Suppose $a_1 \nmid a_2$. Add row 2 to row 1; the first row becomes $(a_1, a_2, 0, \ldots, 0)$. Applying Step 2 replaces $a_1$ by $\gcd(a_1, a_2)$, strictly reducing $\ell(a_1)$. Repeating forces $a_1 \mid a_2$, and similarly $a_1 \mid a_i$ for all $i$. Repeating the argument for $a_2$ in place of $a_1$, and so on, gives $a_i \mid a_{i+1}$ for all $i = 1, \ldots, r-1$.
\end{proof}

\begin{definition}[8.9 — Invariant factors]
The non-zero diagonal entries $a_1, \ldots, a_r$ of the Smith normal form of $A$ are called the \emph{invariant factors} of $A$. They are unique up to unit multipliers, and two $m \times n$ matrices are equivalent if and only if they have the same invariant factors.
\end{definition}

\begin{theorem}[8.10]
Let $A$ be an $m \times n$ matrix over a PID. Then the row rank of $A$ equals the column rank of $A$.
\end{theorem}

\begin{proof}
Choose invertible $P$, $Q$ such that $PAQ$ is in Smith normal form (Theorem~8.8). By Theorem~8.3, multiplying by invertible matrices preserves row and column rank, so
\[
  \mathrm{row\ rank}(A) = \mathrm{row\ rank}(PAQ) = r = \mathrm{col\ rank}(PAQ) = \mathrm{col\ rank}(A),
\]
where $r$ is the number of non-zero diagonal entries.
\end{proof}

\begin{definition}[8.11 — Rank]
The common value of the row rank and column rank of $A$ is called the \emph{rank} of $A$, denoted $\mathrm{rank}(A)$.
\end{definition}

\section*{Lecture 12 — Finitely Generated Modules over a PID}

\begin{theorem}[9.1 — Decomposition Theorem]
Let $R$ be a principal ideal domain, and let $M$ be a finitely generated $R$-module. Then
\[
  M \cong R^s \oplus R/Ra_1 \oplus \cdots \oplus R/Ra_r,
\]
where $s, r \in \mathbb{N} \cup \{0\}$, the $a_i$ are non-zero non-units in $R$, and $a_i \mid a_{i+1}$ for $i = 1, \ldots, r-1$.
\end{theorem}

\begin{proof}
Since $M$ is finitely generated, by Lemma 5.24 there is a surjective $R$-homomorphism $\phi_0 : R^n \to M$, so $M \cong R^n/K$ where $K = \ker \phi_0$. Since $K$ is a submodule of the free module $R^n$, by Lemma 5.25 $K$ is free of rank $m \leq n$. Let $\phi : R^m \xrightarrow{\sim} K$ be an isomorphism and $\{e_1, \ldots, e_m\}$ the standard basis of $R^m$. Writing $\phi(e_j) = (a_{1j}, \ldots, a_{nj})^\top \in R^n$, the map $\phi$ is left-multiplication by the matrix $A = (a_{ij}) \in M_{n \times m}(R)$, so $K = AR^m$.

By Theorem 8.8 (Smith Normal Form), there exist invertible $P \in M_n(R)$ and $Q \in M_m(R)$ such that $PAQ = \mathrm{diag}(a_1, \ldots, a_k, 0, \ldots, 0)$ with $a_1 \mid \cdots \mid a_k$. Since $Q$ is invertible, $\{Qv \mid v \in R^m\} = R^m$, so $AR^m = AQR^m$. Left-multiplication by the invertible $P$ is an automorphism of $R^n$ sending $K = AQR^m$ to $P(AQR^m) = PAQR^m$, so $R^n/K \cong R^n/PAQR^m$. Therefore
\[
  M \cong \frac{R^n}{PAQR^m} \cong \frac{R}{Ra_1} \oplus \cdots \oplus \frac{R}{Ra_k} \oplus R^{n-k},
\]
where the last isomorphism is justified as follows. Left-multiplying $v = (v_1, \ldots, v_m)^\top \in R^m$ by $D = \mathrm{diag}(a_1,\ldots,a_k,0,\ldots,0) \in M_{n \times m}(R)$ gives $(a_1 v_1, \ldots, a_k v_k, 0, \ldots, 0)^\top \in R^n$, and as $v$ ranges over $R^m$ each $a_i v_i$ ranges independently over $Ra_i$. Hence
\[
  PAQR^m = DR^m = Ra_1 \oplus \cdots \oplus Ra_k \oplus \underbrace{0 \oplus \cdots \oplus 0}_{n-k}.
\]
Now apply Theorem 4.33 with $n$ factors, taking $M_i = R$ for every $i$, $A_i = Ra_i$ for $i \leq k$, and $A_i = 0$ for $i > k$:
\[
  \frac{R^n}{PAQR^m} = \frac{R \oplus \cdots \oplus R}{Ra_1 \oplus \cdots \oplus Ra_k \oplus 0 \oplus \cdots \oplus 0} \cong \frac{R}{Ra_1} \oplus \cdots \oplus \frac{R}{Ra_k} \oplus \underbrace{\frac{R}{0} \oplus \cdots \oplus \frac{R}{0}}_{n-k},
\]
and $R/0 \cong R$, so the last $n-k$ summands give $R^{n-k}$. Deleting the zero summands (those where $a_i$ is a unit, giving $R/Ra_i = 0$), renaming the remaining non-unit diagonal entries $a_1, \ldots, a_r$, and setting $s = n - k$ gives
\[
  M \cong \frac{R}{Ra_1} \oplus \cdots \oplus \frac{R}{Ra_r} \oplus R^s. \qedhere
\]
\end{proof}

\begin{definition}[9.2 — Torsion element]
Let $R$ be a principal ideal domain and $M$ an $R$-module. An element $x \in M$ is called a \emph{torsion element} if there exists a non-zero $r \in R$ such that $rx = 0$.
\end{definition}

\begin{definition}[9.4 — Torsion-free element]
Let $R$ be a principal ideal domain and $M$ an $R$-module. A non-zero element $x \in M$ is called a \emph{torsion-free element} if $rx = 0$ for $r \in R$ implies $r = 0$.
\end{definition}

\begin{theorem}[9.6]
Let $R$ be a principal ideal domain and $M$ an $R$-module. Then
\[
  \mathrm{Tor}\, M := \{x \in M \mid x \text{ is a torsion element}\}
\]
is a submodule of $M$.
\end{theorem}

\begin{proof}
We verify the three conditions of Definition 4.8.

\textit{Non-empty.} For any non-zero $r \in R$, $r \cdot 0 = 0$, so $0 \in \mathrm{Tor}\, M$.

\textit{Closed under subtraction.} Let $x, y \in \mathrm{Tor}\, M$, say $rx = 0$ and $sy = 0$ for non-zero $r, s \in R$. Then
\[
  rs(x - y) = s(rx) - r(sy) = 0.
\]
Since $R$ is a PID it is an integral domain, so $rs \neq 0$, hence $x - y \in \mathrm{Tor}\, M$.

\textit{Closed under scalar multiplication.} Let $x \in \mathrm{Tor}\, M$ with $rx = 0$ for some non-zero $r \in R$, and let $a \in R$. Then $r(ax) = a(rx) = 0$, so $ax \in \mathrm{Tor}\, M$.
\end{proof}

\begin{theorem}[9.7]
Let $M$ be a finitely generated module over a principal ideal domain $R$. Then
\[
  M = F \oplus \mathrm{Tor}\, M,
\]
where $F \cong R^s$ for some non-negative integer $s$, and $\mathrm{Tor}\, M \cong R/Ra_1 \oplus \cdots \oplus R/Ra_r$ with the $a_i$ non-zero non-units satisfying $a_1 \mid \cdots \mid a_r$.
\end{theorem}

\begin{proof}
By Theorem 9.1, $M \cong R^s \oplus R/Ra_1 \oplus \cdots \oplus R/Ra_r$. Set $F \cong R^s$ and $T \cong R/Ra_1 \oplus \cdots \oplus R/Ra_r$, so $M = F \oplus T$. We show $T = \mathrm{Tor}\, M$.

\textit{$T \subseteq \mathrm{Tor}\, M$.} We say $c \in R$ \emph{annihilates} a module $N$ if $cy = 0$ for every $y \in N$, i.e.\ multiplying by $c$ sends everything in $N$ to zero. We show that the single non-zero element $a_r$ annihilates $T$, which makes every element of $T$ torsion.

\textit{Step 1: $a_r \in Ra_i$ for each $i$.} Since $a_1 \mid a_2 \mid \cdots \mid a_r$, each $a_i$ divides $a_r$, i.e.\ $a_r = q_i a_i$ for some $q_i \in R$. Hence $a_r \in Ra_i$.

\textit{Step 2: $a_r$ annihilates each summand $R/Ra_i$.} An element of $R/Ra_i$ is a coset $y + Ra_i$ with $y \in R$, and
\[
  a_r(y + Ra_i) = a_r y + Ra_i = y q_i a_i + Ra_i = Ra_i,
\]
since $yq_ia_i \in Ra_i$ ($R$ is commutative). The coset $Ra_i$ is the zero element of $R/Ra_i$, so $a_r(y + Ra_i) = 0$.

\textit{Step 3: $a_r$ annihilates $T$.} An element of $T$ is a tuple $t = (t_1, \ldots, t_r)$ with $t_i \in R/Ra_i$. Scalar multiplication acts coordinatewise, so by Step 2
\[
  a_r t = (a_r t_1, \ldots, a_r t_r) = (0, \ldots, 0) = 0.
\]
Since $a_r \neq 0$, every $t \in T$ is a torsion element, so $T \subseteq \mathrm{Tor}\, M$.

\textit{$\mathrm{Tor}\, M \subseteq T$.} Let $x \in \mathrm{Tor}\, M$ and write $x = x_1 + x_2$ with $x_1 \in F$ and $x_2 \in T$. Since $T \subseteq \mathrm{Tor}\, M$, the element $x_2$ is torsion, and therefore so is $x_1 = x - x_2$. Say $rx_1 = 0$ for some non-zero $r \in R$. Under the isomorphism $F \cong R^s$, $x_1$ corresponds to some $(y_1, \ldots, y_s) \in R^s$, and $r(y_1, \ldots, y_s) = 0$ gives $ry_i = 0$ for each $i$. Since $R$ is an integral domain and $r \neq 0$, each $y_i = 0$, so $x_1 = 0$. Thus $x = x_2 \in T$.
\end{proof}

\begin{theorem}[9.8]
Let $R$ be a principal ideal domain. Then any non-zero ideal $P \neq R$ is prime if and only if it is maximal.
\end{theorem}

\begin{proof}
$(\Leftarrow)$ Every maximal ideal is prime by Theorem 6.21.

$(\Rightarrow)$ Let $P = (a)$ be a non-zero prime ideal. Suppose for a contradiction that $P$ is not maximal, so there exists an ideal $M$ with $P \subsetneq M \subsetneq R$. Since $R$ is a PID, $M = (b)$ for some $b \in R$. Then $(a) \subseteq (b)$ implies $a = bx$ for some $x \in R$. Since $b \notin (a)$ and $(a)$ is prime, the relation $a = bx \in (a)$ forces $x \in (a)$, say $x = ay$ for some $y \in R$. Then $a = bx = bay$, and since $R$ is an integral domain and $a \neq 0$, we may cancel to get $1 = by$. But then $b$ is a unit, so $M = (b) = R$, a contradiction.
\end{proof}

\begin{corollary}
Let $R$ be a principal ideal domain and $P$ a non-zero prime ideal. Then $R/P$ is a field.
\end{corollary}

\begin{proof}
Since $P$ is a non-zero prime ideal in the PID $R$, Theorem 9.8 gives that $P$ is maximal. It remains to show $R/P$ is a field, i.e.\ that every non-zero element $a + P \in R/P$ is invertible. Since $a \notin P$, the ideal $Ra + P$ strictly contains $P$. As $P$ is maximal, $Ra + P = R$, so there exist $r \in R$ and $p \in P$ with $ra + p = 1$. Then $(r + P)(a + P) = ra + P = 1 + P$, so $a + P$ is invertible.
\end{proof}

\begin{theorem}[9.9 — Uniqueness of Decomposition]
Let $R$ be a principal ideal domain and $M$ a finitely generated $R$-module. Suppose
\[
  M \cong R^s \oplus R/Ra_1 \oplus \cdots \oplus R/Ra_u \quad \text{and} \quad M \cong R^t \oplus R/Rb_1 \oplus \cdots \oplus R/Rb_v,
\]
where the $a_i$ and $b_j$ are non-zero non-units with $a_1 \mid \cdots \mid a_u$ and $b_1 \mid \cdots \mid b_v$. Then $s = t$, $u = v$, and $Ra_i = Rb_i$ for all $i$.
\end{theorem}

\begin{proof}
\textit{Step 1: $s = t$.} By Theorem 9.7, both decompositions yield $M = F \oplus \mathrm{Tor}\,M$ with $F \cong R^s$ and $F' \cong R^t$ respectively. Then
\[
  R^s \cong F \cong M/\mathrm{Tor}\,M \cong F' \cong R^t,
\]
so $s = t$ by Theorem 5.21. Let $T = \mathrm{Tor}\,M$; it remains to show $u = v$ and $Ra_i = Rb_i$.

\textit{Step 2: $u = v$.} For an $R$-module $X$ and a prime $p \in R$, define $X_p := \{x \in X \mid px = 0\}$. By the corollary to Theorem 9.8, $R/Rp$ is a field, so $T_p$ is a vector space over $R/Rp$. From the first decomposition of $T$,
\[
  \left(\frac{R}{Ra_i}\right)_p = \begin{cases} R/Rp & \text{if } p \mid a_i, \\ 0 & \text{if } p \nmid a_i, \end{cases}
\]
so $\dim_{R/Rp}(T_p)$ equals the number of $a_i$ divisible by $p$. Now choose a prime $p$ dividing $a_1$; since $a_1 \mid a_i$ for all $i$, $p$ divides every $a_i$, giving $\dim_{R/Rp}(T_p) = u$. Applying the same reasoning to the second decomposition, the dimension also equals the number of $b_j$ divisible by $p$, which is at most $v$, so $u \leq v$. By symmetry (choosing a prime dividing $b_1$), $v \leq u$, hence $u = v$.

\textit{Step 3: $Ra_u = Rb_u$.} Since $a_i \mid a_u$ for all $i$, the element $a_u$ annihilates every summand in the first decomposition, so $a_u T = 0$. Applied to the second decomposition, $a_u \cdot (R/Rb_u) = 0$, which means $a_u \in Rb_u$, so $Ra_u \subseteq Rb_u$. By symmetry $Rb_u \subseteq Ra_u$, hence $Ra_u = Rb_u$.

\textit{The remainder of the proof is non-examinable.} It remains to show $Ra_i = Rb_i$ for $i < u$ by downward induction on $i$. The idea is: assuming $Ra_i = Rb_i$ for $k \leq i \leq u$, one multiplies both decompositions of $T$ by a carefully chosen element $x \in R$ built from a prime $p$ and the exact $p$-adic valuation of $a_u$. This zeros out the first $k-1$ summands in one decomposition but not the other, producing two decompositions of $xT$ with a different number of non-zero summands — contradicting the fact (established by the dimension argument in Step 2) that any two such decompositions must have equally many non-zero summands.
\end{proof}

\begin{theorem}[9.10 — Fundamental Theorem of Finitely Generated Abelian Groups]
Let $A$ be a finitely generated abelian group. Then
\[
  A \cong \mathbb{Z}^s \oplus \mathbb{Z}/a_1\mathbb{Z} \oplus \cdots \oplus \mathbb{Z}/a_r\mathbb{Z},
\]
where $s \geq 0$, the $a_i$ are non-zero non-units in $\mathbb{Z}$, and $a_1 \mid a_2 \mid \cdots \mid a_r$. This decomposition is unique.
\end{theorem}

\begin{proof}
Every abelian group is a $\mathbb{Z}$-module (with $n \cdot a$ meaning $a + \cdots + a$ or its inverse), and $\mathbb{Z}$ is a PID. So Theorem 9.1 applies directly with $R = \mathbb{Z}$, giving a decomposition $A \cong \mathbb{Z}^s \oplus \mathbb{Z}/a_1\mathbb{Z} \oplus \cdots \oplus \mathbb{Z}/a_r\mathbb{Z}$ where each $\mathbb{Z}/a_i\mathbb{Z} = R/Ra_i$ is a cyclic torsion summand. Uniqueness of $s$ and of the ideals $\mathbb{Z}a_i$ (hence the $a_i$ up to sign) then follows from Theorem 9.9.

If $A$ is finite, then $s = 0$: the summand $\mathbb{Z}^s$ contains elements of infinite order whenever $s \geq 1$, but a finite group has no such elements.
\end{proof}

\section*{Lecture 13 — Rational Canonical Form}

\begin{definition}[{The $F[x]$-module structure on $V$}]
Let $F$ be a field, $V$ a finite-dimensional vector space over $F$, and $T \in \mathrm{Hom}_F(V,V)$ a linear map with matrix $A$ relative to some basis. We make $V$ into an $F[x]$-module by defining, for any polynomial $f(x) = a_0 + a_1x + \cdots + a_mx^m \in F[x]$ and any $v \in V$,
\[
  f(x) \cdot v := a_0 v + a_1(Tv) + \cdots + a_m(T^m v).
\]
This extends the $F$-action on $V$ so that $x \cdot v = Tv$. Since any basis $\{v_1, \ldots, v_n\}$ of $V$ over $F$ also generates $V$ as an $F[x]$-module, $V$ is a finitely generated $F[x]$-module. Since $F[x]$ is a PID (as $F$ is a field), the structure theorem of Lecture 12 applies.
\end{definition}

\begin{theorem}[10.1]
Let $T \in \mathrm{Hom}_F(V,V)$, and let $f_1(x), \ldots, f_n(x)$ be the invariant factors of $A - xI$, where $A$ is a matrix of $T$. Then
\[
  V \cong \frac{F[x]}{(f_1(x))} \oplus \cdots \oplus \frac{F[x]}{(f_n(x))}.
\]
\end{theorem}

\noindent\textit{(Proof non-examinable.)} The key steps are: define the surjective $F[x]$-homomorphism $\phi: F[x]^n \to V$ by $e_i \mapsto v_i$, then show $\ker \phi = (A - xI)F[x]^n$, giving $V \cong F[x]^n / (A-xI)F[x]^n$. Since invertible changes of basis preserve this quotient up to isomorphism, one may apply Smith Normal Form (Theorem 8.8) to put $A - xI$ in diagonal form $\mathrm{diag}(f_1(x), \ldots, f_n(x))$. The image of $F[x]^n$ under this diagonal matrix is $(f_1(x)) \oplus \cdots \oplus (f_n(x))$, so Theorem 4.33 with $M_i = F[x]$ and $A_i = (f_i(x))$ for every $i$ gives
\[
  V \cong \frac{F[x] \oplus \cdots \oplus F[x]}{(f_1(x)) \oplus \cdots \oplus (f_n(x))} \cong \frac{F[x]}{(f_1(x))} \oplus \cdots \oplus \frac{F[x]}{(f_n(x))}.
\]

\begin{remark}[10.2]
The polynomials $f_1(x), \ldots, f_n(x)$ are called the \emph{invariant factors} of $A$ (or of $T$). After discarding any unit factors, the remaining non-unit invariant factors are called the invariant factors of the $F[x]$-module $V$.
\end{remark}

\begin{definition}[10.3 — $T$-cyclic space]
Let $T \in \mathrm{Hom}_F(V,V)$. If $V$, regarded as an $F[x]$-module via $T$, is cyclic, then $V$ is called a \emph{$T$-cyclic space} and $T$ is called a \emph{cyclic endomorphism} of $V$.
\end{definition}

\begin{theorem}[10.4]
Let $W$ be a submodule of the $F[x]$-module $V$, and let $T \in \mathrm{Hom}_F(V,V)$ satisfy $TW \subseteq W$. Then $W$ is a $T$-cyclic subspace if and only if there exists $w \in W$ such that $\{w, Tw, \ldots, T^{k-1}w\}$ is a basis of $W$ for some $k \in \mathbb{N}$.
\end{theorem}

\begin{proof}
$(\Rightarrow)$ Suppose $W$ is $T$-cyclic, generated by $w$, so $W = F[x] \cdot w$ in the sense of Definition 4.13. Since $\dim V = n$ is finite, the vectors $w, Tw, T^2w, \ldots, T^nw$ must be linearly dependent over $F$. Let $T^rw$ be the first of these linearly dependent on the preceding ones, so $\{w, Tw, \ldots, T^{r-1}w\}$ are linearly independent and
\[
  T^rw = a_0w + a_1(Tw) + \cdots + a_{r-1}(T^{r-1}w), \quad a_i \in F.
\]
We claim every $T^iw$ is a linear combination of $\{w, Tw, \ldots, T^{r-1}w\}$. This is clear for $i \leq r-1$ and holds for $i = r$ by the equation above. For $i > r$, apply $T^{i-r}$ to the equation to get
\[
  T^iw = a_0(T^{i-r}w) + \cdots + a_{r-1}(T^{i-1}w),
\]
and each term on the right is expressible in terms of $\{w, \ldots, T^{r-1}w\}$ by the induction hypothesis. Since $W = F[x] \cdot w$, every element of $W$ is of the form
\[
  f(x) \cdot w = (b_0 + b_1 x + \cdots + b_m x^m) \cdot w = b_0w + b_1(Tw) + \cdots + b_m(T^mw),
\]
using the $F[x]$-module action from the definition. Since every $T^iw$ lies in $\mathrm{span}\{w, \ldots, T^{r-1}w\}$, this set spans $W$. Combined with linear independence, $\{w, Tw, \ldots, T^{r-1}w\}$ is a basis of $W$ with $k = r$.

$(\Leftarrow)$ Suppose $\{w, Tw, \ldots, T^{k-1}w\}$ is a basis of $W$. Any element $c_0w + c_1(Tw) + \cdots + c_{k-1}(T^{k-1}w)$ of $W$ equals $(c_0 + c_1x + \cdots + c_{k-1}x^{k-1}) \cdot w \in F[x] \cdot w$ by the module action, so $W \subseteq F[x] \cdot w$. Since $w \in W$ and $W$ is closed under the $F[x]$-action, $F[x] \cdot w \subseteq W$. Hence $W = F[x] \cdot w = (w)$, so $W$ is $T$-cyclic by Definition 10.3.
\end{proof}

\begin{definition}[10.5 — Companion matrix]
Let $f(x) = a_0 + a_1x + \cdots + a_{k-1}x^{k-1} + x^k$ be a monic polynomial over $F$. The \emph{companion matrix} of $f(x)$ is the $k \times k$ matrix
\[
C(f) =
\begin{pmatrix}
0      & 0      & \cdots & 0      & -a_0     \\
1      & 0      & \cdots & 0      & -a_1     \\
0      & 1      & \cdots & 0      & -a_2     \\
\vdots & \vdots &        & \vdots & \vdots   \\
0      & 0      & \cdots & 1      & -a_{k-1}
\end{pmatrix}.
\]
\end{definition}

\begin{definition}[10.7 — Minimal polynomial of $T$ on $W$]
Let $W$ be a $T$-cyclic subspace generated by $w$. Since $F[x]$ is a PID, the annihilator $\mathrm{Ann}(w) = \{f(x) \in F[x] \mid f(x) \cdot w = 0\}$ is a principal ideal $(f(x))$ for a unique monic generator $f(x)$. This $f(x)$ is called the \emph{minimal polynomial} of $T$ on $W$.
\end{definition}

\begin{theorem}[10.6]
Let $T \in \mathrm{Hom}_F(V,V)$ and let $W$ be a $T$-cyclic subspace of $V$ with basis $\{w, Tw, \ldots, T^{k-1}w\}$. Then the matrix of $T$ on $W$ is the companion matrix of the minimal polynomial $f(x)$ of $T$ on $W$.
\end{theorem}

\begin{proof}
Since $\{w, Tw, \ldots, T^{k-1}w\}$ is a basis of $W$, the vector $T^kw$ is a linear combination of these basis elements, say
\[
  T^kw = a_0w + a_1(Tw) + \cdots + a_{k-1}(T^{k-1}w), \quad a_i \in F.
\]
In terms of the $F[x]$-module action ($x \cdot v = Tv$), this says $(a_0 + a_1x + \cdots + a_{k-1}x^{k-1} - x^k) \cdot w = 0$, so the polynomial $g(x) := a_0 + a_1x + \cdots + a_{k-1}x^{k-1} - x^k$ lies in $\mathrm{Ann}(w) = (f(x))$.

Since $F[x]$ is a PID, $\mathrm{Ann}(w)$ is a principal ideal $(f(x))$ for some monic $f(x)$ of degree $s$. The inclusion $g(x) \in (f(x))$ gives $s \leq k$. On the other hand, $f(x) \cdot w = 0$ means $T^sw$ is a linear combination of $\{w, Tw, \ldots, T^{s-1}w\}$. Since $\{w, Tw, \ldots, T^{k-1}w\}$ are linearly independent, this forces $s \geq k$. Hence $s = k$.

Since $g(x) \in (f(x))$, write $g(x) = c \cdot f(x)$ for some $c \in F$. Both $g$ and $f$ are monic of degree $k$ with leading coefficient $-1$ and $1$ respectively, so $c = -1$ and $f(x) = -g(x) = x^k - a_{k-1}x^{k-1} - \cdots - a_1x - a_0$.

The matrix of $T$ on $W$ relative to $\{w, Tw, \ldots, T^{k-1}w\}$ has columns given by the images of each basis vector: $T(T^iw) = T^{i+1}w$ for $i < k-1$, contributing $1$s on the subdiagonal, and the last column records the coefficients $a_0, \ldots, a_{k-1}$ of $T^kw$. By Definition 10.5, this is exactly the companion matrix of $f(x)$.
\end{proof}

\section*{Lecture 14}

\begin{theorem}[10.8 — Rational Canonical Form]
Let $V$ be a finite-dimensional vector space over a field $F$ and let $T \in \mathrm{Hom}_F(V,V)$. Then there exists a basis of $V$, as an $F[x]$-module, with respect to which the matrix of $T$ is
\[
\begin{pmatrix}
B_1 & & \\
& \ddots & \\
& & B_r
\end{pmatrix},
\]
where $B_i$ is the companion matrix of a monic polynomial $f_i(x)$, for $i = 1, \ldots, r$, such that $f_1(x) \mid f_2(x) \mid \cdots \mid f_r(x)$. This matrix is called the \emph{rational canonical form} of $T$, and it is unique.
\end{theorem}

\begin{proof}
\textit{Step 1: $V$ is a torsion $F[x]$-module.} Since $\dim_F V = n$, we have $\dim_F \mathrm{Hom}_F(V,V) = n^2$. So the $n^2+1$ elements $I, T, T^2, \ldots, T^{n^2}$ must be linearly dependent over $F$. Let $T^k$ be the first element linearly dependent on the preceding ones, so
\[
  T^k = a_0 I + a_1 T + \cdots + a_{k-1}T^{k-1}, \quad a_i \in F.
\]
Applying both sides to any $v \in V$ and using the $F[x]$-module action $x \cdot v = Tv$ gives
\[
  (a_0 + a_1 x + \cdots + a_{k-1}x^{k-1} - x^k) \cdot v = 0.
\]
Since this polynomial is nonzero, every $v \in V$ is a torsion element (Definition 9.2), so $V = \mathrm{Tor}\,V$.

\textit{Step 2: Decompose $V$ using Theorem 9.1.} $V$ is a finitely generated $F[x]$-module: any $F$-basis of $V$ also generates $V$ as an $F[x]$-module, since $F[x]$ is a PID. Since $V = \mathrm{Tor}\,V$, the free part in Theorem 9.1 vanishes, giving
\[
  V = V_1 \oplus \cdots \oplus V_r,
\]
where each $V_i \cong F[x]/(f_i(x))$ is a cyclic $F[x]$-module and $f_1(x) \mid f_2(x) \mid \cdots \mid f_r(x)$.

\textit{Step 3: Represent $T$ on each $V_i$.} Each $V_i$ is cyclic as an $F[x]$-module, hence a $T$-cyclic subspace (Definition 10.3). By Theorem 10.6, there exists a basis $\mathcal{B}_i$ of $V_i$ with respect to which $T|_{V_i}$ has matrix $B_i$, the companion matrix of $f_i(x)$.

\textit{Step 4: Assemble.} Since $V = V_1 \oplus \cdots \oplus V_r$, the union $\mathcal{B} = \mathcal{B}_1 \cup \cdots \cup \mathcal{B}_r$ is a basis of $V$, and the matrix of $T$ with respect to $\mathcal{B}$ is $\mathrm{diag}(B_1, \ldots, B_r)$.

\textit{Uniqueness.} The uniqueness of the decomposition subject to the divisibility chain, established in Theorem 9.9, implies the polynomials $f_i(x)$ — and hence the companion matrices $B_i$ — are uniquely determined by $T$.
\end{proof}

\begin{definition}[11.1 — Bi-additive map]
Let $M$, $N$, and $P$ be abelian groups. A map $f : M \times N \to P$ is called \emph{bi-additive} if
\[
  f(m_1 + m_2,\, n) = f(m_1, n) + f(m_2, n)
  \quad\text{and}\quad
  f(m,\, n_1 + n_2) = f(m, n_1) + f(m, n_2)
\]
for all $m, m_1, m_2 \in M$ and $n, n_1, n_2 \in N$.
\end{definition}

\begin{definition}[11.3 — Balanced map]
Let $M$ be a right $R$-module and $N$ a left $R$-module. A map $f : M \times N \to P$ is called \emph{balanced} if it is bi-additive and satisfies
\[
  f(mr, n) = f(m, rn)
\]
for all $m \in M$, $n \in N$, and $r \in R$.
\end{definition}

\begin{definition}[11.5 — Tensor product]
A \emph{tensor product} of a right $R$-module $M$ and a left $R$-module $N$ is a pair $(T, \phi)$, where $T$ is an abelian group and $\phi : M \times N \to T$ is a balanced map, satisfying the following universal property: for every abelian group $P$ and every balanced map $f : M \times N \to P$, there exists a unique group homomorphism $\bar{f} : T \to P$ such that $\bar{f} \circ \phi = f$.
\end{definition}

\begin{remark}[11.6 — Uniqueness of the tensor product]
If $(T', \phi')$ is another tensor product of $M$ and $N$, then there exists an isomorphism $\alpha : T \to T'$ such that $\phi' = \alpha \circ \phi$. Thus the tensor product is unique up to isomorphism. We write $M \otimes_R N$ for the tensor product and $m \otimes n$ for $\phi(m, n)$.
\end{remark}

\begin{theorem}[{11.2 —  Existence and construction of $M \otimes_R N$}]
Let $M$ be a right $R$-module and $N$ a left $R$-module. The tensor product $M \otimes_R N$ exists. Concretely, it is realised as $T(M,N) = A(M,N)/B(M,N)$, where $A(M,N)$ is the free abelian group with basis $M \times N$, and $B(M,N)$ is the subgroup generated by all elements of the form
\begin{align*}
  &(m_1 + m_2, n) - (m_1, n) - (m_2, n),\\
  &(m, n_1 + n_2) - (m, n_1) - (m, n_2),\\
  &(mr, n) - (m, rn),
\end{align*}
for all $m, m_1, m_2 \in M$, $n, n_1, n_2 \in N$, and $r \in R$.
\end{theorem}

\noindent\textit{(Proof non-examinable.)} 

\section*{Lecture 15}

\begin{theorem}[11.7]
Let $M$ and $N$ be right and left $R$-modules, respectively. Then for all $m, m_1, m_2 \in M$, $n, n_1, n_2 \in N$, and $r \in R$:
\begin{enumerate}[label=(\roman*)]
  \item $(m_1 + m_2) \otimes n = m_1 \otimes n + m_2 \otimes n$,
  \item $m \otimes (n_1 + n_2) = m \otimes n_1 + m \otimes n_2$,
  \item $mr \otimes n = m \otimes rn$.
\end{enumerate}
\end{theorem}

\noindent\textit{(Proof non-examinable.)} These follow directly from the fact that $B(M,N)$ maps to zero under $\phi$, so the defining relations of $B(M,N)$ become identities in $M \otimes_R N$.

\begin{theorem}[11.8]
Let $M$ and $N$ be right and left $R$-modules, respectively.
\begin{enumerate}[label=(\alph*)]
  \item For $\lambda \in \mathbb{Z}$, $m \in M$, $n \in N$: $(\lambda m) \otimes n = \lambda(m \otimes n) = m \otimes \lambda n$. In particular, $0 \otimes n = 0 = m \otimes 0$ and $(-m) \otimes n = -(m \otimes n) = m \otimes (-n)$.
  \item Every element of $M \otimes_R N$ is a finite sum $\sum_i m_i \otimes n_i$ for $m_i \in M$, $n_i \in N$.
\end{enumerate}
\end{theorem}

\noindent\textit{(Proof non-examinable.)} For (a), fix $n$; the map $m \mapsto m \otimes n$ is a group homomorphism, so it commutes with $\lambda$-multiples. For (b), every element of $A(M,N)$ is a finite $\mathbb{Z}$-linear combination $\sum \lambda_i(m_i, n_i)$, and since $\lambda_i(m_i \otimes n_i) = (\lambda_i m_i) \otimes n_i$ by (a), every element of $M \otimes_R N$ is a finite sum of simple tensors.

\begin{remark}[Module structure on $M \otimes_R N$ — 11.3]
If $R$ is commutative, then $M \otimes_R N$ can be given the structure of a right $R$-module by defining $(m \otimes n)r = mr \otimes n$. This is well defined because $R$ being commutative ensures the map $fr : (m,n) \mapsto mr \otimes n$ vanishes on $B(M,N)$.
\end{remark}

\begin{theorem}[11.9]
Let $M$ be a right $R$-module. Then $M \otimes_R R \cong M$ as additive abelian groups. If $R$ is commutative, then $M \otimes_R R \cong M$ as $R$-modules.
\end{theorem}

\noindent\textit{(Proof non-examinable.)} The map $\bar{f} : M \otimes_R R \to M$ defined by $\bar{f}(m \otimes r) = mr$ is a group isomorphism, with inverse $g : m \mapsto m \otimes 1$.

\begin{theorem}[11.10]
Let $M$, $N$, and $P$ be $R$-modules over a commutative ring $R$. Then
\begin{enumerate}[label=(\roman*)]
  \item $M \otimes_R N \cong N \otimes_R M$ as $R$-modules,
  \item $(M \otimes_R N) \otimes_R P \cong M \otimes_R (N \otimes_R P)$ as $R$-modules.
\end{enumerate}
\end{theorem}

\noindent\textit{(Proof non-examinable.)} For (i), the map $m \otimes n \mapsto n \otimes m$ induces an isomorphism $M \otimes_R N \to N \otimes_R M$; the inverse is constructed symmetrically. Part (ii) is analogous.

\begin{definition}[{Tensor product of homomorphisms — §11.4}]
Let $f : M \to M'$ and $g : N \to N'$ be $R$-homomorphisms. The \emph{tensor product} $f \otimes g : M \otimes_R N \to M' \otimes_R N'$ is the homomorphism defined by
\[
  (f \otimes g)(m \otimes n) = f(m) \otimes g(n).
\]
When $R$ is commutative, $f \otimes g$ is an $R$-homomorphism.
\end{definition}

\begin{theorem}[11.12]
Let $f : M \to M'$, $f' : M' \to M''$, $g : N \to N'$, $g' : N' \to N''$ be $R$-homomorphisms. Then
\[
  (f' \otimes g')(f \otimes g) = f'f \otimes g'g.
\]
In particular, if $f$ and $g$ are isomorphisms then $f \otimes g$ is also an isomorphism.
\end{theorem}

\noindent\textit{(Proof non-examinable.)} Follows directly from the definition: $(f' \otimes g')(f \otimes g)(m \otimes n) = f'f(m) \otimes g'g(n)$.

\begin{theorem}[11.13]
Let $(M_i)_{i \in I}$ and $(N_j)_{j \in J}$ be families of right and left $R$-modules, respectively. Then
\[
  \left(\bigoplus_{i \in I} M_i\right) \otimes_R \left(\bigoplus_{j \in J} N_j\right) \cong \bigoplus_{(i,j) \in I \times J} (M_i \otimes_R N_j),
\]
via the isomorphism $\left(\sum_i x_i\right) \otimes \left(\sum_j y_j\right) \mapsto \sum_{i,j} x_i \otimes y_j$.
\end{theorem}

\noindent\textit{(Proof non-examinable.)} The map $f$ sending $(\sum_i x_i, \sum_j y_j) \mapsto \sum_{i,j} x_i \otimes y_j$ is balanced, inducing a homomorphism $\bar{f}$ in one direction. The inclusions $\lambda_i : M_i \hookrightarrow \bigoplus M_k$ and $\mu_j : N_j \hookrightarrow \bigoplus N_\ell$ give homomorphisms $\lambda_i \otimes \mu_j$, which extend to an inverse $\bar{g}$; then $\bar{f}$ and $\bar{g}$ are mutually inverse isomorphisms.

\begin{theorem}[11.14]
Let $R$ be a commutative ring, $f : M \to M'$ and $g : N \to N'$ be $R$-homomorphisms. If $f$ and $g$ are surjective, then $f \otimes g$ is surjective.
\end{theorem}

\noindent\textit{(Proof non-examinable.)} Every simple tensor $m' \otimes n' \in M' \otimes_R N'$ can be written as $f(m) \otimes g(n) = (f \otimes g)(m \otimes n)$ since $f$ and $g$ are surjective; hence every generator of $M' \otimes_R N'$ is in the image of $f \otimes g$.

\end{document}